\section{The variance identity: savings $D_{\mathrm{var}}$ from fixed marginals}\label{sec:variance}

\paragraph{What this section proves and how to read it.}
When a replacement keeps the same one-half marginals, its energy
against $\Nbr$ equals a diagonal energy minus a product energy (a
variance). The trivial bound keeps only the diagonal ($\le BC_XY$);
this section bounds the product energy below by $gC_XY$, giving the
improved class bound $B-g$ allowed by (CP5). The total saving
$D_{\rm var}$ is subtracted in $\Theta$ as $(63/64)D$ (a
\emph{saving}, called a \emph{credit} in the text; the word
\emph{coupon} below means one proved term in a single coefficient: a subtracted
saving $g$ for price coefficients, or an added energy lower bound
for the mean constant $m$). Four records per
middle block are treated (\emph{eligible} means: replaced at fixed
marginals and annihilating constants). Reading-order note: the
final $D_{\rm var}$ in \cref{cor:variance:final} uses dominator
caps from \cref{sec:envelopes}, and \cref{sec:envelopes} in turn
uses the variance savings for aligned classes; the per-block
inequalities here do not depend on \cref{sec:envelopes}, only the
final assembly does---read the two sections as: per-block
identities first, assembly last. Inputs:
\cref{sec:middle-records}. Output: $D_{\rm var}>0.0099$ (baseline).

The lower bounds are not positive semidefinite orders at $\iota$. They restore
the product energy from the root of the reference in four steps. (A) At the
chamber cut of the fold that produces the record, the actual half law of the
root and the law induced by the record are mutually absolutely continuous,
with a pointwise density bound and a chi-square bound
(\cref{lem:record:root-marginal}; for the positive roots of the late
pieces, \cref{lem:record:root-marginal}(iv)); the chi-square bound rests on the
second moment of a neighbourhood count, which a modular evaluation of a
character sum makes small (\cref{lem:record:weighted-moment}). (B) The root is
a probability coupling at that cut, so its energy is at most that
of its own half diagonal; step (A) compares that diagonal with
the record. (C) The reference is a positive mixture of the root
and of laws of controlled energy, so a stated positive fraction
of $Y$ is recovered from the root energy
(\cref{lem:record:restoration}). (D) At the reserved cut the
record's own Markov operator, in its own actual marginal space,
has centred norm $k_\phi<1$ (\cref{subsec:variance:centred}),
i.e. it is a contraction towards the projection onto constants. The two
matching records are not folds; they are reached from the two fold records by
a polarization at the chamber cut $\vartheta_0$
(\cref{lem:record:paired-e-bridge,lem:record:paired-e-joint}).
\Cref{lem:var-gain} turns the four steps into a gain. For a selected matching
product the exterior bracket adds positive background terms, which a joint
restoration bounds below (\cref{prop:background-variance}). The energies
$Y[\kappa]$, $Y^{\Nbr}[\kappa]$, $Y^{\Nbrz}[\kappa]$ of selected laws, the price
$C_X=\prod_{\Blk\in X}q_\Blk^{2a}$, the diagonal copies $\operatorname{Diag}_\vartheta$,
the product laws $J_\vartheta$ and the relations
\ref{it:record:E1}--\ref{it:record:E3}, \ref{it:record:P1}--\ref{it:record:P4} and
\ref{it:record:P3b} are those of \cref{sec:comparison-lemmas}
(\cref{lem:energy-toolkit}); $Y=Y_\mu$.

\subsection{The four records treated here}\label{subsec:variance:eligible}

The variance saving applies to only four records of a middle
block: the two nonaligned folds $F_1,F_2$ and the two matchings
$E_0,E_1$, called the \emph{eligible} records. Each is replaced
by a law with exactly the same actual half marginals at $\iota$,
so the difference annihilates constants and part of its price can
be recovered as a variance, as the following subsections show.
This subsection fixes the four laws, their actual half marginals,
and the two comparison laws---the diagonal copy and the product
of the marginals---between which the variance is measured.

Let $\Blk$ be a middle block (\cref{def:record:exact-blocks}) with least prime
$b=b_\Blk$, and put $r=r_\Blk=b^{-1/2}$ and $q=r^2$. We write $x\to y$ if
$\chi(y-x)=1$ at every prime of $\Blk$; the laws below are products over the
primes of $\Blk$ of the laws at one prime. We use the components of the
reference of $\Blk$ (\cref{subsec:middle-reference-law}): the \emph{two-layer $F$
law} $\lambda_F$, uniform on six distinct labels
$(y_{o,\ell})_{o\in\{0,1,2\},\,\ell\in\{0,1\}}$ with $y_{0,\ell}\to y_{2,\ell'}$ for
all $\ell,\ell'$, and the \emph{two-layer $E$ law} $\lambda_E$, uniform on six
distinct labels; a map $\phi\colon\Z/6\Z\to\{0,1,2\}$ places the label
$y_{\phi(u),\ell}$ at the six slots $(u,v,\ell)$, $v\in\Z/6\Z$; the type-$\rm D$
components $\Pi_{F_j,\Blk}$ are the images of $\lambda_F$ under the folds $F_j$
of \cref{def:patterns}, of weight $w_{\rm D}r$ each in $\widetilde\mu_\Blk$, and the
type-$\rm S$ components $\Pi_{E_s,\Blk}$ the images of $\lambda_E$ under the matchings
$E_s$, of weight $w_{\rm S}r$ each; and the \emph{root} $\Pi_{\mathrm L,\Blk}$ is the
uniform law on twelve distinct labels $(y_{u,\ell})_{u\in\Z/6\Z,\,\ell\in\{0,1\}}$
with $y_{u,\ell}\to y_{\rot u,\ell'}$ for all $u,\ell,\ell'$, which contains all four
layer lifts of the six directed edges $0\to2\to4\to0$ and $1\to5\to3\to1$; this
is the root of the blocks of the pieces starting below $2^{30}$. On the seven
pieces starting at $2^{30},2^{31},2^{32},2^{34},2^{36},2^{40},2^{45}$ the root is the positive root
$\Gamma_{\Blk,\ulaw}$ of \cref{def:positive-root}, a positive reweighting of the
uniform law with the same support; it is invariant under the
chamber automorphisms used below, the layer interchange and independent
square-affine maps at the primes of $\Blk$
(\cref{lem:positive-root-properties,lem:imp:coherent39-law}),
and every statement of this section about the root holds for it, with the
density and chi-square bounds of its half law given by
\cref{lem:record:root-marginal}(iv) in place of (i)--(iii). For $j\in\{0,1,2\}$ the
\emph{fold chamber} of $F_j$ is the chamber cut $\vartheta_{2j}$ of the first
coordinate, in the notation of \cref{subsec:chamber-geometry}: its reflection
is $u\mapsto4j-1-u$ (that is $5-u$, $3-u$ and $1-u$) and its plus half is
$H_{2j}=\{2j,2j+1,2j+2\}$ (indices mod $6$). By
\cref{lem:record:fold-identification}(a)--(d): the root is invariant under the
chamber automorphisms (for a positive root by
\cref{lem:positive-root-properties}(b) and \cref{lem:imp:coherent39-law}), so at each $\vartheta_{2j}$ it couples
two copies of one half law $\nu_{\mathrm L,\Blk}$, which, read in the order
$(2j,2j+1,2j+2)$, lives on the support of $\lambda_F$ and does not depend on $j$;
$\Pi_{F_j,\Blk}=\operatorname{Diag}_{\vartheta_{2j}}(\lambda_F^{(j)})$, where
$\lambda_F^{(j)}$ places the labels of the object $o$ at the object $2j+o$; and
at $\iota$ the records $\Pi_{F_j,\Blk}$, $\Pi_{E_s,\Blk}$ are symmetric, with
actual half marginal the layer-zero triple law of $\lambda_F$ placed by $F_j$,
respectively the uniform law on ordered triples of distinct labels placed by
$E_s$.

\begin{definition}[Eligible records and repairs]\label{def:record:eligible}
The \emph{eligible records} of a middle block $\Blk$ are
$\phi\in\{F_1,F_2,E_0,E_1\}$. For such $\phi$ let $\Pi_{\phi,\Blk}$ be the
normalized reference record, $m_{\phi,\Blk}$ its actual half marginal at $\iota$
(the same on both halves), $D_{\phi,\Blk}=\operatorname{Diag}_\iota(m_{\phi,\Blk})$
its \emph{selected diagonal}, and $J_{\phi,\Blk}$ the law under which the two
$\iota$-halves are independent, each with law $m_{\phi,\Blk}$ (the
\emph{product of the actual half marginals}). A \emph{repair} of $\phi$ is a
signed measure $\Pi'_{\phi,\Blk}$ of mass one on the labels of $\Blk$ that is a
finite real combination of nonnegative measures, each invariant under
independent square-affine maps at the primes of $\Blk$ and with the actual half
marginal $m_{\phi,\Blk}$ on both sides of $\iota$. The \emph{record difference}
$\Delta_{\phi,\Blk}$ is the operator on $L^2(m_{\phi,\Blk})$ with
\[
 \ip{\psi}{\Delta_{\phi,\Blk}\psi'}_{m_{\phi,\Blk}}
 =\int\psi(z_+)\psi'(z_-)\,d(\Pi'_{\phi,\Blk}-\Pi_{\phi,\Blk})(z),
\]
and $\Delta^{\rm s}_{\phi,\Blk}$, $\Delta^{\rm a}_{\phi,\Blk}$ are its symmetric and
skew parts. In the comparison law of the middle range the repair of $F_j$ is
the signed whole-block fold record of \cref{def:record:mixture}(c),(d), a
combination of three positive product laws, each with both half marginals
$m_{F_j,\Blk}$ (\cref{lem:signed-components}), and the repair of $E_s$ is the
exact-marginal matching law of \cref{lem:record:e-law}, tensorized over the
primes of $\Blk$ (\cref{lem:record:exact-marginal-e}).
\end{definition}

Both $\Pi_{\phi,\Blk}$ and its repair have the actual half marginal
$m_{\phi,\Blk}$ on both sides of $\iota$, so
\begin{equation}\label{eq:record:zero-marginal}
 \Delta_{\phi,\Blk}\one=\Delta_{\phi,\Blk}^*\one=0,\qquad
 \Delta_{\phi,\Blk}=(I-\Jproj)\Delta_{\phi,\Blk}(I-\Jproj),
\end{equation}
where $\Jproj$ is the projection onto the constants of $L^2(m_{\phi,\Blk})$; the
same holds for $\Delta^{\rm s}_{\phi,\Blk}$ and $\Delta^{\rm a}_{\phi,\Blk}$. The
record of weight $w_\phi\in\{w_{\rm D}r,w_{\rm S}r\}$ contributes
$w_\phi(\Pi'_{\phi,\Blk}-\Pi_{\phi,\Blk})$, normalized by
$\widetilde\mu_\Blk(1)\ge1$, to $\nu_\Blk-\lambda_\Blk$: the first stage changes
only the one-layer folded records (\cref{def:record:frame}), so the two-layer
records enter $\lambda_\Blk$ and $\mu_\Blk$ with the same weight. In the
per-record decomposition of \cref{def:record:comparison-law}(b) its norm
envelopes are $\tilde\pi^{\rm s}D_{\phi,\Blk}$ and $\tilde\pi^{\rm a}D_{\phi,\Blk}$
for any $\tilde\pi^{\rm s}\ge w_\phi\|\Delta^{\rm s}_{\phi,\Blk}\|$ and
$\tilde\pi^{\rm a}\ge w_\phi\|\Delta^{\rm a}_{\phi,\Blk}\|$
(\cref{rem:cp:charging}(ii)), with prices $\tilde\pi^{\rm s}q_\Blk^{2a}$ and
$\tilde\pi^{\rm a}q_\Blk^{2a}$; the norm bounds are those of
\cref{prop:signed-bounds} for $F_j$ and of \cref{lem:record:exact-marginal-e} for
$E_s$. The two records of one kind in a block have differences of equal
symmetric and skew norms: that of $F_2$ is $UD^*U$, with $D$ that of $F_1$ and
$U$ a self-adjoint unitary, so that its symmetric and skew parts are
$UD^{\rm s}U$ and $-UD^{\rm a}U$ (\cref{def:record:mixture}(d),
\cref{prop:signed-bounds}(e)), and $E_0$, $E_1$ place one and the same law. So
the two fold records of a block carry the same prices, and so do its two
matching records. Every positive law constructed in this section
(the records, their repairs, their diagonals and their products)
is a positive measure,
invariant under independent translations at the primes of its blocks, whose
normalization has uniform one-slot marginals and does not depend on the
indices (O1); the only operator-norm bounds at a cut applied to
these laws, the contraction ($d=1$) of a probability coupling
and the centred norms of
a record's own Markov operator at $\iota$, condition on no shared label (O2). So
the laws are admissible (\cref{def:admissible}), and
\ref{it:record:E1}--\ref{it:record:E3} apply to them.

The comparison at the fold chamber of $F_j$ is the following statement.

\begin{hypothesis}[Root density, (A)]\label{hyp:variance:A}
For $j\in\{1,2\}$, at the fold chamber $\vartheta_{2j}$ of $F_j$ (reflection
$u\mapsto3-u$ for $F_1$, $u\mapsto1-u$ for $F_2$), the actual half law
$\nu_{\mathrm L,\Blk}$ of the root on the objects $H_{2j}$ and the law
$\lambda_F^{(j)}$, whose $\vartheta_{2j}$-diagonal copy is the record
$\Pi_{F_j,\Blk}$, satisfy $M_\Blk^{-1}\le d\nu_{\mathrm L,\Blk}/d\lambda_F^{(j)}
\le M_\Blk$.
\end{hypothesis}

\Cref{lem:record:root-marginal} proves hypothesis (A) with an explicit
$M_\Blk$, for the uniform root in (i)--(iii) and for the positive roots of the
late pieces in (iv), and also bounds the chi-square divergence of the two laws.
The record compared with the root is the two-layer record itself, which the
first stage does not change, not its reserved-cut marginal. The
comparison is made at a chamber cut because no comparison is available at
$\iota$: there the half laws of the folded records are singular to that of the
root, and no positive semidefinite order at $\iota$ bounds the product energy
from below (\cref{rem:record:variance-n1}).

\subsection{Comparing the principal half law with the record half law}\label{subsec:variance:root}

Restoration needs the chi-square distance between two half laws
at a chamber cut: the half law of the principal (root) record and
the half law of an eligible record. At one prime this distance is
governed by the second moment of a signed neighbourhood count---
how many common neighbours, with prescribed signs, four labels
have. That second moment is computed first; one of the graph
moments in it is evaluated through the Ahlgren--Ono identity
(\cref{thm:ahlgren-ono}), the others by Weil's bound. At one prime $p\equiv3\pmod4$ we use $q=1/p$, $r=p^{-1/2}$ and $k_p=(p-3)/4$,
the exact number of common successors, and of common predecessors, of two
distinct elements of $\Fp$ (\cref{lem:paley-triple}(a)). Signed
neighbourhoods of four distinct labels and their soft counts $\mathfrak S$ are
those of \cref{lem:record:signed-nbhd}: with $\upsilon_+=11r+6q$ and
$\upsilon_-=11r+70q$, a signed neighbourhood has between
$\frac p{16}(1-\upsilon_-)$ and $\frac p{16}(1+\upsilon_+)$ elements, and at
least $\frac p{16}(1-\upsilon_\times)$ of them avoid two further labels, where
\[
 \upsilon_\times=\upsilon_-+32q=11r+102q .
\]

\paragraph{The second moment of a neighbourhood count.}
The chi-square comparison needs the second moment of the fluctuation of a
neighbourhood count. For independent uniform $x=(x_1,x_2,x_3,x_4)\in\Fp^4$ and
the signs $s=(-1,-1,1,1)$ put $s_A=\prod_{i\in A}s_i$ and
\[
 \mathcal W=\sum_{\substack{A\subseteq\{1,2,3,4\}\\|A|\in\{3,4\}}}s_A
 \sum_{z\in\Fp}\prod_{i\in A}\chi(z-x_i),\qquad
 P_{\rm adm}=\prod_{i\in\{1,2\},\ k\in\{3,4\}}\bigl(1+\chi(x_k-x_i)\bigr).
\]
For a finite simple graph $G$ on $\{0,\dots,n-1\}$ let
$M_G=\E\prod_{\{u,v\}\in G,\,u<v}\chi(\xi_v-\xi_u)$ for independent uniform
$\xi_0,\dots,\xi_{n-1}\in\Fp$. Let $G_5$ be the graph on five vertices with the
eight edges $01$, $02$, $03$, $04$, $13$, $14$, $23$, $24$ (the complete
tripartite graph with parts $\{0\}$, $\{1,2\}$, $\{3,4\}$); $G_{12}$ the complete
tripartite graph on six vertices with parts $\{0,1\}$, $\{2,3\}$, $\{4,5\}$; and
$G_{10}$ the graph on six vertices with the ten edges $01$, $03$, $04$, $05$,
$12$, $14$, $15$, $23$, $25$, $34$.

\begin{lemma}[Second moment of the neighbourhood count]
\label{lem:record:weighted-moment}
Let $p\equiv3\pmod4$ be prime.
\begin{enumerate}[label=(\alph*)]
\item The moment satisfies
 \begin{equation}\label{eq:record:weighted-identity}
  \frac{\E[P_{\rm adm}\mathcal W^2]}{p^2}=(4+8q)M_{G_5}+M_{G_{12}}+24M_{G_{10}}
  +5q-4q^2+14q^3-15q^4,
 \end{equation}
 $|M_{G_5}|,|M_{G_{12}}|\le q$ and $|M_{G_{10}}|\le qr$. Consequently
 $\E[P_{\rm adm}\mathcal W^2]/p^2\le10q+24qr+4q^2+14q^3$.
\item With $T_\lambda=\sum_{y\in\Fp}\chi(y(y-1)(y-\lambda))$,
 \begin{equation}\label{eq:record:modular-moment}
  M_{G_5}=\frac{p-1}{p^4}\sum_{\lambda\in\Fp}\chi(\lambda)T_\lambda^2,\qquad
  -(1-q)(2qr+q^2)\le M_{G_5}\le(1-q)(2qr-q^2),
 \end{equation}
 and $M_{G_{12}}=(1-5q)M_{G_5}+q^2(1-q)(1-2q)$. Consequently
 \[
  \frac{\E[P_{\rm adm}\mathcal W^2]}{p^2}=(5+3q)M_{G_5}+24M_{G_{10}}+5q-3q^2
  +11q^3-13q^4\le5q+34qr+11q^3 .
 \]
\end{enumerate}
\end{lemma}

\begin{proof}
\emph{The identity.} Expand $P_{\rm adm}$ into its $16$ products over subsets
of the four cross pairs, and $\mathcal W^2$ into $25$ products over pairs
$(A,A')$, with independent summation labels $z,z'$. Dividing by $p^2$, each of
the $400$ terms is $\pm$ a moment $\E\prod\chi(\xi_v-\xi_u)$ of an oriented
simple graph on the six vertices $x_1,x_2,x_3,x_4,z,z'$, with $z,z'$ uniform;
all these graphs are subgraphs of the complete tripartite graph with parts
$\{x_1,x_2\},\{x_3,x_4\},\{z,z'\}$. The moments are reduced by four exact rules,
valid for every prime $p\equiv3\pmod4$:
(G1) a vertex of degree one makes the moment zero, since
$\sum_\xi\chi(\xi-\eta)=0$; (G2) an odd number of edges makes it zero, since
$\xi\mapsto-\xi$ preserves the uniform law and multiplies the product by
$\chi(-1)^{\#\text{edges}}$; (G3) relabelling the vertices changes the moment
by the sign $\chi(-1)$ per reversed edge, and an automorphism reversing an odd
number of edges forces the moment to vanish; (G4) for a vertex $v$ of degree
two with neighbours $u,w$, $\E_{\xi_v}\chi(\xi_v-\xi_u)\chi(\xi_v-\xi_w)=
\one[\xi_u=\xi_w]-q$ (\cref{lem:paley-triple}(b)), which removes $v$ and adds the
graph with $u$ and $w$ identified; a repeated edge is
$\chi(t)^2=1-\one[t=0]$ and is expanded in the same way, and a loop gives zero.
Every step lowers the number of vertices or of vertices of degree two, so the
reduction terminates in polynomials in $q$ and moments of graphs with all
degrees at least three. Among the graphs that occur these are $G_5$, $G_{12}$
and $G_{10}$, and collecting the $400$ reductions gives
\eqref{eq:record:weighted-identity}. The bookkeeping of the $400$ terms is the
finite symbolic computation of \cref{cert:record:weighted-identity}.

\emph{The bounds for $M_{G_5}$ and $M_{G_{12}}$.} Suppose a graph $G$ has two
disjoint edges $u_1v_1,u_2v_2$ whose four endpoints do not span a complete
graph $K_4$. If one of $u_1v_2,u_2v_1$ is missing, let the two sides be
$\{u_1,u_2\}$ and $\{v_1,v_2\}$; otherwise one of $u_1u_2,v_1v_2$ is missing,
and the sides are $\{u_1,v_2\}$ and $\{v_1,u_2\}$. In either case at most three
of the four pairs across the sides are edges, and both selected edges cross.
Fix all other labels. Edges inside a side, and edges to fixed labels, are
multipliers of modulus at most one on the functions of that side; the crossing
edges are the kernel of an operator $T_{\mathcal K}$ of \cref{lem:matching:kernel}
with $k=2$, whose diagonal pairs are the two selected edges and which has at
most one further pair, so that its digraph is acyclic and
$\|T_{\mathcal K}\|\le r^2=q$ by part (a) of that lemma. The fixed-label moment
is therefore at most $q$ in absolute value, and so is $M_G$. In $G_5$ the edges
$01$ and $23$ qualify ($12$ is missing), and in $G_{12}$ the edges $02$ and $34$
($23$ is missing).

\emph{The bound for $M_{G_{10}}$.} Translate $\xi_0=0$. The factor $\chi(\xi_1)$
vanishes at $\xi_1=0$; for $\xi_1=c\ne0$ substitute $\xi_v=c\eta_v$ ($v\ge2$):
each of the ten edge factors acquires $\chi(c)$, and $\chi(c)^{10}=1$. With
$\eta_2=\beta$, $\eta_3=\gamma$,
\[
 M_{G_{10}}=\frac{p-1}{p^5}\sum_{\beta,\gamma}\chi(\gamma)\chi(\beta-1)
 \chi(\gamma-\beta)\,T_\beta T_\gamma,\qquad
 T_\beta=\sum_y\chi(y)\chi(y-1)\chi(y-\beta),
\]
the sums over $\eta_4$ and $\eta_5$ giving $T_\gamma$ and $T_\beta$. The matrix
$H_{\gamma\beta}=\chi(\gamma-\beta)$ satisfies $H^{\mathsf T}H=pI-\Jall$ and
$H\one=0$, so $\|H\|\le\sqrt p$; and $T=-Hv$ with $v_y=\chi(y)\chi(y-1)$,
$\|v\|^2=p-2$, so $\|T\|\le p$. The double sum is $\psi'^{\mathsf T}H\psi$ with
$\psi_\beta=\chi(\beta-1)T_\beta$ and $\psi'_\gamma=\chi(\gamma)T_\gamma$, of norms
at most $p$; hence it is at most $p^{5/2}$ and
$|M_{G_{10}}|\le(p-1)p^{-5/2}\le p^{-3/2}=qr$.

The last assertion of (a) follows from \eqref{eq:record:weighted-identity}, the
three bounds and $-15q^4\le0$.

\emph{The modular evaluation.} (b) Translate $\xi_0=0$ in $M_{G_5}$. The factor
$\chi(\xi_1)$ vanishes at $\xi_1=0$; for $\xi_1=c\ne0$ substitute $\xi_v=c\eta_v$
($v=2,3,4$): each of the eight edge factors acquires $\chi(c)$, and
$\chi(c)^8=1$. With $\lambda=\eta_2$ the edges at the vertex $3$ give
$\chi(\eta_3)\chi(\eta_3-1)\chi(\eta_3-\lambda)$, and those at the vertex $4$ the
same factor in $\eta_4$, so the sums over $\eta_3$ and $\eta_4$ are both
$T_\lambda$, and the factor $\chi(\eta_2)$ of the edge $02$ is $\chi(\lambda)$; this
is the identity of \eqref{eq:record:modular-moment}. By \cref{thm:ahlgren-ono},
$\sum_\lambda\chi(\lambda)T_\lambda^2=-p-a(p)$ with $|a(p)|\le2p^{3/2}$; multiplying
by $(p-1)/p^4$ gives the two bounds, since $(p-1)p^{-3}=(1-q)q^2$ and
$2(p-1)p^{-5/2}=2(1-q)qr$.

\emph{The inversion identity.} Translate $\xi_0=0$ in $M_{G_{12}}$, so that
$M_{G_{12}}=p^{-5}\sum_{\xi_1,\dots,\xi_5}\Phi(\xi)$, where $\Phi$ is the product of
the twelve edge factors at $\xi_0=0$. The four edges at the vertex $0$ give the
factor $\chi(\xi_2)\chi(\xi_3)\chi(\xi_4)\chi(\xi_5)$.

On the part $\xi_1\ne0$, $\Phi$ vanishes unless all five labels are nonzero;
substitute $\xi_v=1/\eta_v$ ($1\le v\le5$). Then $\chi(\xi_v)=\chi(\eta_v)$, and for
an edge $uw$ with $u<w$ between nonzero labels
$\chi(\xi_w-\xi_u)=\chi\bigl((\eta_u-\eta_w)/(\eta_u\eta_w)\bigr)
=-\chi(\eta_w-\eta_u)\chi(\eta_u)\chi(\eta_w)$. Every vertex has degree four in
$G_{12}$, and each edge at a vertex $v\ge1$ contributes one factor $\chi(\eta_v)$,
so these factors give $\chi(\eta_v)^4=1$; the eight edges between the vertices
$1,\dots,5$ give the sign $(-1)^8=1$. What remains is the product of the factors
$\chi(\eta_w-\eta_u)$ over the edges $12,13,14,15,24,25,34,35$, which form a copy
of $G_5$ (centre $1$, parts $\{2,3\}$ and $\{4,5\}$), summed over nonzero labels.
By inclusion and exclusion over the set of vanishing labels: the unrestricted
sum is $p^5M_{G_5}$; prescribing one vanishing label leaves, by translation,
$p^4M_{G_5}$; two adjacent labels cannot vanish together; the two
nonadjacent pairs $\{2,3\}$ and $\{4,5\}$ each contribute
$\sum_{\eta_1\ne0}\bigl(\sum_{y\ne0}\chi(y-\eta_1)\bigr)^2=p-1$, since
$\sum_{y\ne0}\chi(y-\eta_1)=-\chi(-\eta_1)$; and no three labels are pairwise
nonadjacent. So this part contributes $(p^5-5p^4)M_{G_5}+2(p-1)$.

On the part $\xi_1=0$, the four edges at the vertex $1$ give the same factor as
those at $0$, so the product is the indicator that $\xi_2,\dots,\xi_5$ are
nonzero times the four factors of the complete bipartite graph between
$\{2,3\}$ and $\{4,5\}$. For nonzero $a,b$ put
$c(a,b)=\sum_{y\ne0}\chi(y-a)\chi(y-b)$; then $c(a,a)=p-2$ and
$c(a,b)=-1-\chi(ab)$ for $a\ne b$, and this part is
\[
 \sum_{a,b\ne0}c(a,b)^2=(p-1)(p-2)^2+\sum_{a\ne b}\bigl(2+2\chi(ab)\bigr)
 =(p-1)(p^2-2p-2),
\]
as $\sum_{a\ne b}\chi(ab)=-(p-1)$. Adding the two parts and dividing by $p^5$
gives $M_{G_{12}}=(1-5q)M_{G_5}+(p-1)(p-2)/p^4$, which is the inversion identity.

Substituting it in \eqref{eq:record:weighted-identity} gives the displayed
identity of (b). By \eqref{eq:record:modular-moment}, $5+3q>0$ and
$|M_{G_{10}}|\le qr$, its right side is at most $5q+34qr+11q^3$ minus
$8q^2-2q^3+10q^4+4q^2r+6q^3r$, which is nonnegative.
\end{proof}

\begin{remark}[The crude second moment]\label{rem:variance:crude}
The identity matters only for its constant. Under independent uniform labels
only the diagonal pairs $A=A'$ survive in $\E[\mathcal W^2]$, because
$\E_x\chi(z-x)=0$, and $\E_x\chi(z-x)\chi(z'-x)=\one[z=z']-q$; hence
$\E[\mathcal W^2]=\sum_{|A|\in\{3,4\}}[p(1-q)^{|A|}+p(p-1)(-q)^{|A|}]\le5p$, and
$0\le P_{\rm adm}\le16$ gives $\E[P_{\rm adm}\mathcal W^2]/p^2\le80q$. This
already shows that the chi-square divergence of \cref{lem:record:root-marginal}
is $O(q)$. The identity lowers the constant from $80$ to about $10$, and the
modular evaluation of \cref{lem:record:weighted-moment}(b) to about $5$; the
certified credits use the latter.
\end{remark}

\paragraph{The first moment of the four-root count.}
The cut L20 of the uniform root (\cref{subsec:middle-l20}) needs the first
moment of the same signed count, with the four labels distinct. Write
$\mathcal W=\mathcal W_3+\mathcal W_4$, where $\mathcal W_k$ is the part of
$\mathcal W$ with $|A|=k$, and let $D_4$ be the indicator that $x_1,\dots,x_4$
are distinct.

\begin{lemma}[The first moment of the four-root count]\label{lem:l20-star-mean}
Let $p\equiv3\pmod4$ be prime. Then
\begin{equation}\label{eq:l20-star-identities}
 \frac{\E[P_{\rm adm}D_4\mathcal W_3]}p=0,\qquad
 \frac{\E[P_{\rm adm}D_4\mathcal W_4]}p=M_{G_5}+8q^2-33q^3+25q^4 .
\end{equation}
Consequently, let $p>7$, so that $Z_4=(1-q)(1-3q)(1-7q)>0$, and let $\mathsf K$
be the soft kernel of the cut L20 of \cref{subsec:middle-l20}, the product over
the sixteen cross edges of the factors $1+(D_{20})_{ij}\chi(x'_j-x_i)$ with the
signed cross matrix $D_{20}$ of \eqref{eq:l20-cross-matrix}, for independent
halves $x,x'$ with the law $\lambda_{F,p}$. In the expansion of $\E\mathsf K$ over
the sets of cross edges, the sets that form a star (all their edges at one
label, the centre) have total mean of modulus at most
\[
 m_\star=\frac{4\{(1-q)(q^2+2qr)+8q^2+33q^3+25q^4\}}{Z_4}.
\]
Hence, if the nonstar part $T_{\rm ns}$ of the soft kernel and the collision
correction $T_{\rm coll}$ of \cref{subsec:middle-l20} have norms at most
$N_{\rm ns}$ and $D_{\rm col}$, and $\eps_Z=N_{\rm ns}+D_{\rm col}+m_\star<1$, then
the normalizer $Z_0=\E\mathsf K-\ip\one{T_{\rm coll}\one}$ of the L20 cut operator
satisfies $|Z_0-1|\le\eps_Z$, and the normalizer centred at the half-law mean may use
$R_Z=(1-\eps_Z)^{-1}\ge Z_0^{-1}$.
\end{lemma}

\begin{proof}
\emph{The identities.} Expand $D_4$ over the partitions $\pi$ of the four labels,
with the M\"obius coefficients of the partition lattice:
\begin{equation}\label{eq:partition-coefficient}
 D_4=\sum_\pi\mu(\pi)\,\one_\pi,\qquad
 \mu(\pi)=\prod_{C\in\pi}(-1)^{|C|-1}(|C|-1)!\,,
\end{equation}
where $\one_\pi$ is the indicator that $x$ is constant on every block of $\pi$;
expand $P_{\rm adm}$ into its $16$ products over the sets of cross pairs and
$\mathcal W_k$ into its $k$-sets $A$ with the signs $s_A$. Identifying the labels
of each block of $\pi$ leaves $|\pi|$ independent uniform labels and costs the
factor $q^{4-|\pi|}$. After division by $p$, each of the $16\cdot5\cdot15=1200$
terms is therefore $s_A\mu(\pi)q^{4-|\pi|}$ times the moment of the oriented graph
on the merged labels and the summation label $z$, with the chosen cross pairs and
the edges from $z$ to the labels of $A$; edges reversed by the identification
contribute their signs $\chi(-1)=-1$, and a loop makes the term zero. Each moment
is reduced by the rules (G1)--(G4) of the proof of
\cref{lem:record:weighted-moment}, which are exact for every prime $p\equiv3
\pmod4$. The reduced terms with $|A|=3$ cancel, and those with $|A|=4$ add up to
$M_{G_5}+8q^2-33q^3+25q^4$; the bookkeeping of the $1200$ reductions is the
finite symbolic computation of \cref{cert:l20-star-mean}.

\emph{The star means.} Under $\lambda_F\otimes\lambda_F$ the centre $z$ of a star
is a coordinate of one half, uniform on $\Fp$ and independent of the other half,
which carries the leaves; the leaves are distinct labels. By the orientation of
the factors, a star centred in the first half has the sign $(-1)^{\#\text{leaves}}$
in addition to the product of its entries of $D_{20}$. A one-leaf star has mean
zero, as $\E_z\chi(z-y)=0$. A two-leaf star has mean $-q$ times the product of
its two entries, because $\E_z\chi(z-y)\chi(z-y')=-q$ for $y\ne y'$
(\cref{lem:paley-triple}(b)). The rows of $D_{20}$, and also its columns, are of
two kinds (\eqref{eq:l20-cross-matrix}): four have two nonzero entries, of equal
sign, at the labels $b_0,b_1$ of the other half; two have four nonzero entries,
at the labels $a_0,a_1,c_0,c_1$ of the other half, with the signs
$\pm(-1,-1,1,1)$. The products of the entries
over the two-element subsets of the support therefore sum, over the rows and
over the columns, to $4\cdot1+2\cdot(-2)=0$, and the two-leaf means cancel.
Stars with three or four leaves occur only at the four centres of the second
kind, two on each side, and their leaves are the four labels $(a_0,a_1,c_0,c_1)$
of the other half. Under $\lambda_F$ these labels are uniform on the admissible
ordered quadruples (distinct, with $a_i\to c_j$), since each admissible
quadruple has exactly $(p-4)(p-5)$ completions by $(b_0,b_1)$; for independent
uniform labels $P_{\rm adm}D_4$ is $16$ times the indicator of the admissible
event, which has probability $p(p-1)k_p(k_p-1)/p^4=Z_4/16$. Hence
$\E_{\lambda_F}[f]=\E[P_{\rm adm}D_4f]/Z_4$ for every function $f$ of the four labels.
At such a centre the three-leaf stars together have the mean
$\pm\E_{\lambda_F}[\mathcal W_3]/p$, which is zero by
\eqref{eq:l20-star-identities}, and the four-leaf star, whose sign product is
$+1$, has the mean $\E_{\lambda_F}[\mathcal W_4]/p
=(M_{G_5}+8q^2-33q^3+25q^4)/Z_4$. By \eqref{eq:record:modular-moment},
$|M_{G_5}|\le(1-q)(q^2+2qr)$, so the four centres give at most $m_\star$.

The empty set of edges gives one and the nonstar sets give
$\ip\one{T_{\rm ns}\one}$, of modulus at most $N_{\rm ns}$; the collision
correction $\ip\one{T_{\rm coll}\one}$ has modulus at most $D_{\rm col}$. So
$|Z_0-1|\le m_\star+N_{\rm ns}+D_{\rm col}=\eps_Z$, and $Z_0\ge1-\eps_Z>0$.
\end{proof}

With $q=r^2$, $m_\star/r^2$ is a positive increasing function of $r$ for
$0<r<7^{-1/2}$, its numerator being $4(2r+9r^2-2r^3+32r^4+25r^6)$ and its
denominator $Z_4$ decreasing; so the value of $m_\star$ at a threshold bounds it
at every larger prime.

\begin{certificate}[The four-root first moment]\label{cert:l20-star-mean}
The $1200$ terms of the expansion in the proof of
\cref{lem:l20-star-mean} add up to
\eqref{eq:l20-star-identities}, and the signed two-leaf row and
column sums of the cross matrix vanish; both are checked in exact
arithmetic as follows. The script
\path{research/continuation_2026_10_24/upper_arithmetic/l20_star_mean.py}
writes the record
\path{research/continuation_2026_10_24/upper_arithmetic/l20_star_mean.json}
(status \texttt{PASS\_EXACT\_L20\_STAR\_MEAN\_IDENTITY}): it lists the $1200$ terms
of the expansion in the proof of \cref{lem:l20-star-mean}, with the coefficient
\eqref{eq:partition-coefficient}, the sign and the integer reduction of each by
the rules (G1)--(G4), and checks that they add up to
\eqref{eq:l20-star-identities}; it also checks that the signed two-leaf sums of
the rows and of the columns of the cross matrix vanish, and the identities by
direct finite sums at $p=7,11,19$. The record and the script are pinned by
SHA-256 in \path{research/continuation_2026_10_24/verification/selected.json}.
Its check \texttt{l20\_mean} (status
\texttt{PASS\_INDEPENDENT\_L20\_VARIANCE\_MEAN\_SPECTRAL\_GATES}) runs
\path{research/continuation_2026_10_24/upper_arithmetic/review_middle_l20_mean.py},
which asserts the status of the record, its $1200$ terms, the vanishing of the
part with $|A|=3$ and the two vanishing two-leaf sums; its full-tier check
\texttt{fresh\_l20\_mean} reruns
\path{research/continuation_2026_10_24/upper_arithmetic/middle_l20_variance_mean.py},
which asserts the status of the record as well. The continuation record is an
authenticated input of the later ones: \path{research/advance_2026_10_24/verify.py}
asserts its status and re-hashes the files it pins, and the later verifiers authenticate this chain through
\path{research/push_2026_10_25/verify.py} (\cref{cert:closing}).
\end{certificate}

\begin{lemma}[Root completion count]\label{lem:record:root-marginal}
Let $p\equiv3\pmod4$ be prime and put
\begin{gather*}
 \mathfrak a_p=\frac{4q}{1-3q}\cdot\frac{(1+\upsilon_+)^2}
 {(1-\upsilon_-)(1-\upsilon_--16q)},\qquad
 R_p=\frac{(1+\upsilon_+)^2}{(1-\upsilon_\times)(1-\upsilon_\times-16q)
 (1-24\mathfrak a_p)},\\
 Z_4=(1-q)(1-3q)(1-7q),\qquad
 \Delta_2=\Bigl(\frac{5q+34qr+11q^3}{Z_4}\Bigr)^{1/2}+30q,\\
 \Delta_{\max}=11r+30q,\qquad
 D_{\rm full}=(2+\Delta_{\max}+16q)\Delta_2+16q+24\mathfrak a_p(1+\Delta_{\max})^2,\\
 L_{\rm low}=(1-\Delta_{\max})(1-\Delta_{\max}-16q)(1-24\mathfrak a_p),
\end{gather*}
and assume that the \emph{lower factors} $1-\upsilon_\times-16q$ and
$1-24\mathfrak a_p$ are positive. Then $L_{\rm low}>0$ as well, since
$\Delta_{\max}<\upsilon_\times$.
\begin{enumerate}[label=(\roman*)]
\item The half law $\nu_{\mathrm L,p}$ of the root at one prime satisfies
 $R_p^{-1}\le d\nu_{\mathrm L,p}/d\lambda_{F,p}\le R_p$; moreover the
 supremum of this density is at most $R_p$ times its infimum.
\item $\chi_p:=\chi^2(\nu_{\mathrm L,p}\,\|\,\lambda_{F,p})\le
 \min\bigl((R_p-1)^2/R_p,\,D_{\rm full}^2/L_{\rm low}^2\bigr)$.
\item Let $\Blk$ be a middle block whose primes are at least $X$, with
 $\log q_\Blk\le L$, and let $n\ge\lfloor L/\log X\rfloor$ be an integer.
 Then hypothesis (A) holds with $M_\Blk=\prod_{p\in\Blk}R_p\le\bar R(X)^n$,
 and
 \[
  \chi^2(\nu_{\mathrm L,\Blk}\,\|\,\lambda_{F,\Blk})=\prod_{p\in\Blk}(1+\chi_p)-1
  \le\bigl(1+\bar\chi(X)\bigr)^n-1,
 \]
 where $\bar R(X)$ and $\bar\chi(X)$ are the right sides of (i) and (ii)
 evaluated at $q=1/X$ and any fixed $r\ge X^{-1/2}$ for which the two lower
 factors are positive (the notation suppresses this $r$;
 \eqref{eq:variance:pieces} takes $r=1/\lfloor\sqrt X\rfloor$).
\item \emph{Positive roots.} Let $\Blk$ be a middle block of one of the pieces
 starting at $X\in\{2^{30},2^{31},2^{32},2^{34},2^{36},2^{40},2^{45}\}$, whose root is the positive root
 $\Gamma_{\Blk,\ulaw}$ of \cref{def:positive-root}, with $\log q_\Blk\le L$ and
 $n\ge\lfloor L/\log X\rfloor$, and let $R_\ulaw(X)$ and $d_\ulaw(X)^2$ be the
 range ratio and the variance bound of the half marginal of $\Gamma_{p,\ulaw}$ at
 the chamber cut $\vartheta_0$ given by
\cref{prop:root-chamber,imp:env:coherent-parameter,imp:env:coherent-single-pin} at the endpoint
 $X$, whose gates hold there (\cref{cert:positive-root-caps}). Then hypothesis
 (A) holds for the half law $\nu_{\mathrm L,\Blk}$ of $\Gamma_{\Blk,\ulaw}$ with
 $M_\Blk\le R_\ulaw(X)^n$, and
 \[
  \chi^2(\nu_{\mathrm L,\Blk}\,\|\,\lambda_{F,\Blk})\le
  \bigl(1+d_\ulaw(X)^2\bigr)^n-1 .
 \]
\end{enumerate}
\end{lemma}

\begin{proof}
Work at $\vartheta_0$, with half $H_0=\{0,1,2\}$; the half law does not depend
on the chamber. A half configuration is $y=(y_{u,\ell})_{u\in\{0,1,2\},\,\ell\in
\{0,1\}}$, six distinct labels with $y_{0,\ell}\to y_{2,\ell'}$. It completes to a
root configuration by labels $y_{u,\ell}$, $u\in\{3,4,5\}$, subject to
$y_{2,\ell}\to y_{4,\ell'}\to y_{0,\ell''}$ (even triangle) and
$y_{1,\ell}\to y_{5,\ell'}\to y_{3,\ell''}\to y_{1,\ell'''}$ (odd triangle), for all
layer indices, all twelve labels distinct.

\emph{Odd component.} Fix $(y_{1,0},y_{1,1})$ and let $\mathcal O$ be the set
of $(y_{5,0},y_{5,1},y_{3,0},y_{3,1})$ satisfying the odd constraints with the six
odd labels distinct, and $\mathcal Z=|\mathcal O|$. The labels $y_{5,\ell}$ are an
ordered distinct pair of common successors of $y_{1,0},y_{1,1}$ ($k_p(k_p-1)$
choices), and given them the labels $y_{3,\ell}$ are an ordered distinct pair in
the signed neighbourhood of $(y_{5,0},y_{5,1},y_{1,0},y_{1,1})$ with the signs
$(1,1,-1,-1)$. By \cref{lem:record:signed-nbhd},
$\mathcal Z\ge k_p(k_p-1)(p/16)^2(1-\upsilon_-)(1-\upsilon_--16q)$. The number of
elements of $\mathcal O$ with a prescribed value of $y_{5,0}$ is at most
$(k_p-1)(p/16)^2(1+\upsilon_+)^2$; the same bound holds for $y_{5,1}$, and,
choosing $y_{3,0},y_{3,1}$ first as common predecessors of $y_{1,0},y_{1,1}$, for
$y_{3,0}$ and $y_{3,1}$. So under the uniform law on $\mathcal O$ each of the four
free odd labels takes any given value with probability at most $\mathfrak a_p$,
since $1/k_p=4q/(1-3q)$. Finally $\mathcal Z$ does not depend on
$(y_{1,0},y_{1,1})$: the maps $z\mapsto cz+c'$ with $\chi(c)=1$ preserve all the
constraints and act transitively on the ordered distinct pairs with a fixed
value of $\chi(y_{1,1}-y_{1,0})$; $\mathcal O$ depends only on the set
$\{y_{1,0},y_{1,1}\}$, because all four lifts of each edge are present; and the
two orders of the pair have opposite values of $\chi(y_{1,1}-y_{1,0})$, since
$\chi(-1)=-1$.

\emph{Even component.} Given $y$, the pair $(y_{4,0},y_{4,1})$ is an ordered
distinct pair in the signed neighbourhood of $(y_{0,0},y_{0,1},y_{2,0},y_{2,1})$
with the signs $(-1,-1,1,1)$, avoiding $y_{1,0},y_{1,1}$. By
\cref{lem:record:signed-nbhd} the number of such pairs lies between
$(p/16)^2(1-\upsilon_\times)(1-\upsilon_\times-16q)$ and $(p/16)^2(1+\upsilon_+)^2$.
For each such pair, the odd
completions that avoid the six even labels form a subset of $\mathcal O$ of
relative size between $1-24\mathfrak a_p$ and $1$, by the union bound over four
free odd labels and six forbidden values. Hence the number $\mathfrak N(y)$ of
completions of $y$ satisfies
\begin{equation}\label{eq:record:root-count}
 (1-\upsilon_\times)(1-\upsilon_\times-16q)(1-24\mathfrak a_p)
 \le\frac{\mathfrak N(y)}{\mathcal Z(p/16)^2}\le(1+\upsilon_+)^2 .
\end{equation}

(i) The density of $\nu_{\mathrm L,p}$ with respect to the uniform law
$\lambda_{F,p}$ on the half configurations is $\mathfrak N(y)/
\E_{\lambda_F}\mathfrak N$. By \eqref{eq:record:root-count} the values of
$\mathfrak N$ lie in an interval $[c,R_pc]$, $c>0$ ($\mathcal Z$ cancels), and so
does their mean; hence the supremum of the density is at most $R_p$ times its
infimum, and the density lies in $[R_p^{-1},R_p]$.

(ii) A density $g$ with $\E g=1$ and values in $[R^{-1},R]$ satisfies
$\E(g-R^{-1})(R-g)\ge0$, that is, $\E g^2\le R+R^{-1}-1$; this is the first
bound. For the second, under $\lambda_{F,p}$ the four labels
$x=(y_{0,0},y_{0,1},y_{2,0},y_{2,1})$ are uniform on the \emph{admissible}
ordered quadruples (distinct, with $y_{0,\ell}\to y_{2,\ell'}$), since each has
exactly $(p-4)(p-5)$ choices of $(y_{1,0},y_{1,1})$; for independent uniform
labels the admissible event has probability $p(p-1)k_p(k_p-1)/p^4=Z_4/16$. For
admissible $x$ the soft count with the signs $(-1,-1,1,1)$ equals
$p+2+\mathcal W$: the singleton sums vanish, each quadratic sum is $-1$
by \eqref{eq:two-root}, and $\sum_{i<k}s_is_k=-2$. At each of the four labels of
$x$ one factor of the soft count vanishes (at $z=y_{0,0}$, for instance,
$1+\chi(y_{0,0}-y_{2,0})=1-\chi(y_{2,0}-y_{0,0})=0$), so the soft count is exactly
$16$ times the size of the signed neighbourhood. Removing $y_{1,0},y_{1,1}$
deletes $n_1\in\{0,1,2\}$ points, so the number of admissible labels for
$y_{4,\ell}$ is $n=\bar n(1+\beta)$, with $\bar n=p/16$ and
$\beta=\mathcal W/p+(2-16n_1)q$. By \cref{thm:weil},
$|\mathcal W|\le4\cdot2\sqrt p+3\sqrt p$, so $|\beta|\le\Delta_{\max}$. Since
$P_{\rm adm}\ge0$ everywhere and $P_{\rm adm}=16$ on the admissible event,
\cref{lem:record:weighted-moment}(b) gives
\[
 \E_{\lambda_F}\Bigl(\frac{\mathcal W}p\Bigr)^2\le
 \frac{\E[P_{\rm adm}\mathcal W^2]}{p^2Z_4}
 \le\frac{5q+34qr+11q^3}{Z_4},
\]
hence $\|\beta\|_{L^2(\lambda_F)}\le\Delta_2$. Let $\Xi=n(n-1)/\bar n^2$, the
normalized number of ordered distinct pairs for object $4$. Then
$\Xi-1=\beta(2+\beta)-16q(1+\beta)$, so
$\|\Xi-1\|_{L^2(\lambda_F)}\le(2+\Delta_{\max}+16q)\Delta_2+16q$, and
$(1-\Delta_{\max})(1-\Delta_{\max}-16q)\le\Xi\le(1+\Delta_{\max})^2$, both factors on
the left being positive. By the
even-component step, $\varsigma:=\mathfrak N/(\mathcal Z\bar n^2)$ satisfies
$(1-24\mathfrak a_p)\Xi\le\varsigma\le\Xi$. Therefore
$\|\varsigma-1\|_{L^2(\lambda_F)}\le D_{\rm full}$ and $\E\varsigma\ge L_{\rm low}$.
The density is $\varsigma/\E\varsigma$, and
$\chi_p=\Var(\varsigma)/(\E\varsigma)^2\le\E(\varsigma-1)^2/(\E\varsigma)^2\le
D_{\rm full}^2/L_{\rm low}^2$.

(iii) At a block the root and $\lambda_F$ are products over the primes of
$\Blk$, so the density is the product of the prime densities, and chi-square
divergences of product laws satisfy
$1+\chi^2(\bigotimes_p\nu_p\|\bigotimes_p\lambda_p)=\prod_p(1+\chi^2(\nu_p\|
\lambda_p))$. Each of $\upsilon_\pm$, $\upsilon_\times$, $\mathfrak a_p$, $R_p$,
$\Delta_2$, $\Delta_{\max}$, $D_{\rm full}$ and $1/L_{\rm low}$ is a nondecreasing
function of $(q,r)$ on the domain where the two lower factors are
positive (\cref{lem:majorant}); so for $p\ge X$ the hypotheses of (i) and (ii)
hold, $R_p\le\bar R(X)$ and $\chi_p\le\bar\chi(X)$; and $\Blk$ has at most
$\lfloor L/\log X\rfloor\le n$ primes,
since their product is at most $e^L$ and each is at least $X$. As
$\bar R(X)\ge1$, this proves (iii).

(iv) Each selected tilt is invariant under the chamber reflections
$u\mapsto3-u$, $u\mapsto1-u$ and their product $u\mapsto u+2$.
For the pair laws these maps permute the three scores
(\cref{lem:positive-root-properties}(b)); for the coherent law they
permute its physical support dictionary with its coefficients
(\cref{lem:imp:coherent39-law}). Thus the reweighted root retains the
same three chamber identifications. The
translation maps the plus half $H_0$ of $\vartheta_0$ onto the plus half $H_2$ of
$\vartheta_2$, and $H_2$ onto the plus half $H_4$ of $\vartheta_4$, in order, and it
maps $\lambda_F=\lambda_F^{(0)}$ to $\lambda_F^{(1)}$ and $\lambda_F^{(1)}$ to
$\lambda_F^{(2)}$; each reflection exchanges the two halves of its chamber. Hence,
as for the uniform root, the positive root couples at each fold chamber two
copies of one half law, and its density with respect to $\lambda_F^{(j)}$ is the
same function, read in the order $(2j,2j+1,2j+2)$, at $\vartheta_0$, $\vartheta_2$
and $\vartheta_4$. At $\vartheta_0$ and one prime $p\ge X$ this density is $T_\ulaw\one$
for the cut operator $T_\ulaw$ of $\Gamma_{p,\ulaw}$ on $L^2(\lambda_{F,p})$
(\cref{lem:spec:cutkernel}(f)), and
\cref{prop:root-chamber,imp:env:coherent-parameter,imp:env:coherent-single-pin}
give $\sup T_\ulaw\one/\inf T_\ulaw\one\le R_\ulaw(X)$ and
$\|T_\ulaw\one-\one\|^2\le d_\ulaw(X)^2$; these are the conclusions of
\cref{lem:tilt-marginals} for the cut kernel of the uniform root, with the range
ratio $R=R_p$ of (i) and $\delta_0=(\chi^+_p)^{1/2}$, where
$\chi^+_p=D_{\rm full}^2/L_{\rm low}^2$ is the second bound of (ii). As
$T_\ulaw\one$ has mean one,
$\inf T_\ulaw\one\le1\le\sup T_\ulaw\one$, so $R_\ulaw(X)^{-1}\le T_\ulaw\one\le
R_\ulaw(X)$. On the block the positive root and $\lambda_F$ are products over the
primes, and the product clause of \cref{lem:tilt-marginals} gives the range
ratio $\prod_{p\in\Blk}R_{\ulaw,p}\le R_\ulaw(X)^n$ and the chi-square
$\prod_{p\in\Blk}(1+d_{\ulaw,p}^2)-1\le(1+d_\ulaw(X)^2)^n-1$, with
$R_{\ulaw,p}$, $d^2_{\ulaw,p}$ the values at $p$, since $\Blk$ has at most $n$ primes
and $R_\ulaw(X)\ge1$ (proof of (iii)). The block density has mean one, so, as
at one prime, its range ratio gives hypothesis (A) with
$M_\Blk\le R_\ulaw(X)^n$.
\end{proof}

\subsection{Recovering product energy on a common chamber}\label{subsec:variance:restoration}

The variance saving needs a lower bound for the energy of a
product over several blocks. At one common chamber cut this energy
can be recovered (``restored'') from the diagonal of $\lambda$; the
lemma makes the recovery quantitative, with the loss measured by
the nonroot mass $h_\Blk$ of each block.

\begin{lemma}[Restoration on a common chamber]\label{lem:record:restoration}
Let $j\in\{1,2\}$ and let $X$ be a nonempty finite set of middle blocks. Let
$h_X\ge\prod_{\Blk\in X}(1+h_\Blk)-1$, where $h_\Blk=\widetilde\mu_\Blk(1)-1$
is the nonroot mass of $\widetilde\mu_\Blk$,
\begin{equation}\label{eq:record:nonroot-mass}
 h_\Blk=(3w_{\rm D}+2w_{\rm S}+v_{\rm L})r+(6w_{\rm R}+3v_{\rm D}+2v_{\rm S})r^2
 +(2w_{\rm C}+6v_{\rm R})r^3+2v_{\rm C}r^4,
 \qquad r=r_\Blk .
\end{equation}
Let $M_X\ge\prod_{\Blk\in X}M_\Blk$ with $M_\Blk$ as in \cref{hyp:variance:A},
$\chi_X\ge\chi^2(\bigotimes_\Blk\nu_{\mathrm L,\Blk}\|\bigotimes_\Blk
\lambda_{F,\Blk})$ and $C\ge C_X$, and put $l=1-(C-1)h_X$. Assume $0<l\le C$.
Then, with $\Pi_X=\bigotimes_{\Blk\in X}\Pi_{F_j,\Blk}$, $Y[\Pi_X]\ge0$ and
\[
 Y[\Pi_X]\ge(l_0-o(1))Y,\qquad
 l_0=\max\Bigl(\frac l{M_X},\;l-\frac C2\sqrt{\chi_X},\;
 \frac{2l^2}{2l+C\chi_X+\sqrt{\chi_X(4l(C-l)+C^2\chi_X)}}\Bigr).
\]
\end{lemma}

\begin{proof}
All laws are on the labels of $X$; write $\Pi^{\rm rt}=\bigotimes_\Blk
\Pi_{\mathrm L,\Blk}$, $\nu=\bigotimes_\Blk\nu_{\mathrm L,\Blk}$ and
$\lambda=\bigotimes_\Blk\lambda_F^{(j)}$, so that
$\operatorname{Diag}_{\vartheta_{2j}}(\lambda)=\Pi_X$. The fold chamber
$\vartheta_{2j}$ is a cut of the family, so \ref{it:record:E2},
\ref{it:record:P1} and \ref{it:record:P2} apply there with $\mu$ outside. The
measure $\Pi_X$ is positive and diagonal at $\vartheta_{2j}$, so $Y[\Pi_X]\ge0$
by \ref{it:record:E3} and \ref{it:record:E2}.

\emph{(C) The root carries the reference.} Write $\widetilde\mu_\Blk=
\Pi_{\mathrm L,\Blk}+\mathcal M_\Blk$, where $\mathcal M_\Blk\ge0$ is the sum of
the other weighted components of \cref{lem:front12:reference}. Its mass is
$h_\Blk$, which is \eqref{eq:record:nonroot-mass}: the five types
$\rm L,D,S,R,C$ have multiplicities $1,3,2,6,2$ and degrees $0,1,1,2,3$, a
two-layer component of type $t$ has weight $w_tr^{\deg t}$ and a one-layer
diagonal weight $v_tr^{1+\deg t}$. Expanding $\bigotimes_{\Blk\in X}\mu_\Blk=
\prod_\Blk(1+h_\Blk)^{-1}\bigotimes_\Blk(\Pi_{\mathrm L,\Blk}+\mathcal M_\Blk)$,
every term other than $\Pi^{\rm rt}$ is a positive measure whose normalization
is a product of components, with uniform one-slot marginals, and their total
mass is $h=\prod_\Blk(1+h_\Blk)-1\le h_X$. Since $Y[\bigotimes_\Blk\mu_\Blk]=Y$,
\ref{it:record:E1} gives $Y[\Pi^{\rm rt}]\ge(1+h-(1+o(1))hC_X)Y\ge(l-o(1))Y$.

\emph{(B) Contraction and comparison.} The root $\Pi^{\rm rt}$ is a probability
coupling at $\vartheta_{2j}$ with both half marginals $\nu$, so the contraction
\ref{it:record:P1} gives $Y[\Pi^{\rm rt}]\le
Y[\operatorname{Diag}_{\vartheta_{2j}}(\nu)]$; no positivity of $\Pi^{\rm rt}$ at
the cut is used. By hypothesis (A), $\nu\le M_X\lambda$, so the density
domination \ref{it:record:P2} gives $Y[\operatorname{Diag}_{\vartheta_{2j}}(\nu)]\le
M_XY[\Pi_X]$. Thus $Y[\Pi_X]\ge(l/M_X-o(1))Y$.

\emph{One-sided variation.} The laws $\nu$ and $\lambda$ are invariant under
independent translations at each prime of $X$, and so are the two parts of the
Jordan decomposition $\nu-\lambda=\kappa_+-\kappa_-$, whose masses are
$\frac12\|\nu-\lambda\|_{\rm var}\le\frac12\sqrt{\chi_X}$ by Cauchy--Schwarz. The
normalized diagonal copies of $\kappa_\pm$ have uniform one-slot marginals, and
$Y[\operatorname{Diag}_{\vartheta_{2j}}(\kappa_-)]\ge0$ by \ref{it:record:E2} and
\ref{it:record:E3}. Hence
$Y[\operatorname{Diag}_{\vartheta_{2j}}(\nu)]-Y[\Pi_X]\le
Y[\operatorname{Diag}_{\vartheta_{2j}}(\kappa_+)]\le(1+o(1))\frac12\sqrt{\chi_X}\,C_XY$,
and $Y[\Pi_X]\ge(l-\frac C2\sqrt{\chi_X}-o(1))Y$.

\emph{Bounded variance.} For a half-label tuple $z$ let $o_z$ be the uniform law
on its translates $z+c$, where $c=(c_p)_p$ adds one $c_p\in\Fp$ to all labels at
each prime $p$ of $X$, and put $\eta(z)=Y[\operatorname{Diag}_{\vartheta_{2j}}
(o_z)]/Y$. The diagonal copy of $o_z$ is a positive normalized law with uniform
one-slot marginals, fixed samplewise by the reflection of $\vartheta_{2j}$, so
$0\le\eta(z)\le C':=(1+o(1))C_X$. Translation invariance gives
$\nu=\sum_z\nu(z)o_z$ and $\lambda=\sum_z\lambda(z)o_z$, hence
$Y[\operatorname{Diag}_{\vartheta_{2j}}(\nu)]=Y\,\E_\nu\eta$ and $Y[\Pi_X]=Y\,\bar\eta$
with $\bar\eta=\E_\lambda\eta$. With $u=d\nu/d\lambda-1$,
\[
 \E_\nu\eta-\bar\eta=\E_\lambda[u(\eta-\bar\eta)]
 \le\bigl(\chi_X\Var_\lambda(\eta)\bigr)^{1/2}
 \le\bigl(\chi_X\bar\eta(C'-\bar\eta)\bigr)^{1/2},
\]
because $\eta^2\le C'\eta$. Since $\E_\nu\eta\ge Y[\Pi^{\rm rt}]/Y\ge l-o(1)$ by
(C) and (B), and $C'\le C+o(1)$, we get
$l-o(1)\le\bar\eta+(\chi_X\bar\eta(C-\bar\eta))^{1/2}$. If $\bar\eta<l$, squaring
gives $(1+\chi_X)\bar\eta^2-(2l+C\chi_X)\bar\eta+l^2\le o(1)$; the discriminant of
the left side is $\chi_X(4l(C-l)+C^2\chi_X)\ge0$, and its smaller root,
rationalized, is the third expression in $l_0$, which is at most $l$. So
$\bar\eta$ is at least that expression minus $o(1)$ in either case.
\end{proof}

Write $\mathcal F(C,h,M,\chi)$ for the largest of $0$ and the three expressions
of $l_0$ formed with $l=1-(C-1)h$ if $l>0$, and $\mathcal F(C,h,M,\chi)=0$ if
$l\le0$. As $l\le1\le C$ always and $Y[\Pi_X]\ge0$, the lemma gives
$Y[\Pi_X]\ge(\mathcal F(C,h_X,M_X,\chi_X)-o(1))Y$ under its other hypotheses.
In the construction the cap is a product of piece caps: if every block of the
piece $i$ has $\log q_\Blk\le L_i$ and $a_*\ge a$, then every block of $X$ lying
in the piece $i$ has $q_\Blk^{2a}\le C_i:=e^{2a_*L_i}$, so $C_X$ is at most the
product $\prod_{\Blk\in X}C_{i(\Blk)}$ of the caps of the pieces of its blocks
(\cref{lem:piece-caps}), and the lemma is applied with this product as $C$; the
construction takes $a_*=a$ and the numbers $L_i$ of \cref{lem:piece-caps}.

\begin{remark}[The three branches]\label{rem:variance:branches}
The second expression in $l_0$ never exceeds the third: the smaller root
$\bar\eta_*$ satisfies $l-\bar\eta_*=(\chi_X\bar\eta_*(C-\bar\eta_*))^{1/2}\le
\frac C2\sqrt{\chi_X}$. It is kept because the certificate evaluates all three
expressions, and zero (\cref{cert:variance:fresh}); the certificate also checks
the condition $0\le l\le C/2$, whereas the lemma needs only $l\le C$ (at the
final caps $l\le1<C/2$ holds automatically, every piece cap being larger than
$2$). The pointwise expression $l/M_X$ needs neither the chi-square bound nor the
variance arguments, and it shows the mechanism, but the final restorations rest
on the chi-square expressions: at the final caps the bounded-variance expression
is the largest at every piece and every pair of pieces (at the first piece about $0.98104$,
against about $0.98075$ for the one-sided and $0.71382$ for the pointwise
expression), and on the seven pieces with positive roots $M_i>10^7$.
The bounded-variance expressions therefore provide the useful restoration
on these pieces; the same exact formulas are reevaluated for their own
marginals in \cref{cert:variance:fresh}.
\end{remark}

\subsection{Centred norms at the reserved cut}\label{subsec:variance:centred}

The last step converts a lower bound for the energy of a product of eligible
records into one for the product of its actual half marginals, by the centred
interpolation \ref{it:record:P3} at $\iota$. It needs, for each record, the norm
on the mean-zero space of its own Markov operator at $\iota$, in its own actual
marginal space: for the fold records the centred norm $\kappa_F(q,r)$ of
\cref{lem:record:old-f-marginal}(c), an admissible function of $(q,r)$ that the
certificate evaluates at the piece start, and for the matching records the norm
$3/(b-3)$ of the disjoint-triple operator (\cref{lem:record:paired-e-bridge}(b)
below). For a product of records the centred norm is the largest of the prime
centred norms, not their sum (\cref{lem:sector}(c),
\cref{lem:record:old-f-marginal}(d)), because each prime operator fixes the
constants of its own actual marginal on both sides; we write $\kappa_F(\Blk)$
for the largest value of $\kappa_F(1/p,p^{-1/2})$ at the primes $p$ of $\Blk$.
Inside $\kappa_F$, the kernel of two disjoint cross
edges whose four endpoints do not span a complete rectangle has norm at most
$q$ for independent uniform labels, which is case (a) of
\cref{lem:matching:kernel}, and \cref{lem:record:projection} transfers the bound
to the actual marginal (\cref{lem:old-f-law}(b)).

\subsection{Variance savings for $E_0,E_1$ via polarization}\label{subsec:variance:matching}

The records $E_0,E_1$ are matchings, not folds, and no chamber diagonal is a
matching record. Their energies are reached from those of the two fold records
by a polarization at the chamber cut $\vartheta_0$. For a middle block $\Blk$ let
$G_{e,\Blk}$ and $G_{o,\Blk}$ be the images of $\lambda_E$ under the merges
$G_e=(1,0,0,0,0,1)$ and $G_o=(0,0,1,1,0,0)$ of \cref{def:patterns} (two objects
used); they are positive normalized laws with uniform one-slot marginals,
diagonal at $\vartheta_0$. By the corner table of
\cref{lem:record:fold-identification}(f) (\cref{tab:corner}), at $\vartheta_0$
the cut matrices of $\Pi_{E_0,\Blk}$ and $G_{e,\Blk}$ live on the corner
$\mathcal U$ of the plus-half configurations $(\beta,y,y)$, those of
$\Pi_{E_1,\Blk}$ and $G_{o,\Blk}$ on the corner $\mathcal V$ of the configurations
$(y,y,\beta)$, and $\Pi_{F_1,\Blk}$, $\Pi_{F_2,\Blk}$ connect the two corners; here
$\beta$ is the label pair of the object that is not repeated. Let $f_2,h_2$ be
the two-layer row coefficients of the piece of $\Blk$
(\cref{lem:front12:reference}(i); the index $2$ counts the layers), put
$\varepsilon_\Blk=(f_2+h_2)r_\Blk$, and
$\kappa_E(\Blk)=3/(b_\Blk-3)$. For a finite set $X$ of middle blocks write
$E_{\mathbf s}=\bigotimes_{\Blk\in X}\Pi_{E_{s_\Blk},\Blk}$, $\mathbf s\in\{0,1\}^X$,
for the $2^{|X|}$ \emph{matching embeddings},
$J_{\mathbf s}=\bigotimes_\Blk J_{E_{s_\Blk},\Blk}$ for their products of actual
half marginals, and $E^X_0$, $E^X_1$ for the two \emph{aligned} embeddings, with
$E_0$, respectively $E_1$, on every block.

\begin{lemma}[Paired $E$ bridge]\label{lem:record:paired-e-bridge}
Let $\Blk$ be a middle block.
\begin{enumerate}[label=(\alph*)]
\item As cut matrices at $\vartheta_0$ of laws on the labels of $\Blk$,
 \[
  \Pi_{E_0,\Blk}+\Pi_{E_1,\Blk}+\varepsilon_\Blk(G_{e,\Blk}+G_{o,\Blk})
  \succeq_{\vartheta_0}\Pi_{F_1,\Blk}+\Pi_{F_2,\Blk}.
 \]
\item At $\iota$ the Markov operator $T_E$ of $\Pi_{E_s,\Blk}$ on
 $L^2(m_{E_s,\Blk})$ is self-adjoint and $\|T_E-\Jproj\|\le\kappa_E(\Blk)$. More
 generally, every matching embedding $E_{\mathbf s}$ of a finite set $X$ of
 middle blocks is symmetric at $\iota$, and its Markov operator on $L^2$ of its
 actual half marginal has centred norm at most $\max_{\Blk\in X}\kappa_E(\Blk)$.
\end{enumerate}
\end{lemma}

\begin{proof}
(a) As in the proof of \cref{lem:front12:reference}(iii), the two-layer laws
of $\Blk$ form a five-type datum (\cref{def:five-type-datum}) on the label
pairs of one object: the middle law $\nu$ and the conditional references
$\nu^\xi$ of \cref{lem:front12:reference}(i), the matching law $\lambda_E$, its
two-object marginal, the fold law $\lambda_F$ and the root. By
\cref{lem:front12:reference}(i) the middle object has the law $\nu$ under
$\lambda_F$ as well as under $\lambda_E$ (this uses $p\equiv3\pmod4$), and
$\|T^\xi_{\rm F}-\Jproj\|\le f_2r_\Blk$, $\|T^\xi_{\rm E}-\Jproj\|\le h_2r_\Blk$
on $L^2(\nu^\xi)$ for every $\xi$. The components of this datum at
$\vartheta_0$ are $\Pi_{F_j,\Blk}$, $\Pi_{E_s,\Blk}$, $G_{e,\Blk}$ and $G_{o,\Blk}$,
read on the slot array through the first coordinate. So (a) is
\cref{lem:tr:polarization}(a) with $e_{\rm F}=f_2r_\Blk$, $e_{\rm E}=h_2r_\Blk$ and
$\varepsilon=\varepsilon_\Blk$: the form of the difference of the two sides is at
least the average over $\xi$ of $(u_\xi-v_\xi)^2\ge0$, with $u_\xi,v_\xi$ the
conditional means of the two corner functions. As $\nu^\xi$ gives no mass to
$\xi$, the datum also satisfies the hypothesis of \cref{lem:tr:polarization}(b).

(b) By \cref{lem:record:fold-identification}(d), $\Pi_{E_s,\Blk}$ is symmetric
at $\iota$ and its actual half marginal is uniform on ordered triples of
distinct labels; given the layer-zero triple, the layer-one triple is uniform
among those disjoint from it. So $T_E$ is the tensor product over the primes of
$\Blk$ of the disjoint-triple operators of \cref{sec:frontier:empty-triple-negative},
which are self-adjoint, and part (c) of that lemma gives
$\|T_E-\Jproj\|\le3/(b_\Blk-3)$. An embedding $E_{\mathbf s}$ is a product of such
records, hence symmetric at $\iota$, and its operator is the tensor product of
the operators of its blocks, each fixing the constants of its actual marginal
on both sides; so \ref{it:record:P4} bounds its centred norm by the largest
block value (\cref{lem:sector}(c)).
\end{proof}

\begin{lemma}[Matching selections]\label{lem:record:paired-e-joint}
Let $X$ be a set of $k\ge1$ middle blocks, with $\varepsilon_\Blk$, $G_{e,\Blk}$,
$G_{o,\Blk}$ as above, $G_{0,\Blk}=G_{e,\Blk}$, $G_{1,\Blk}=G_{o,\Blk}$ and
$G_\Blk=G_{e,\Blk}+G_{o,\Blk}$.
\begin{enumerate}[label=(\alph*)]
\item At the chamber cut $\vartheta_0$,
 \[
  \bigotimes_{\Blk\in X}\bigl(\Pi_{E_0,\Blk}+\varepsilon_\Blk G_{e,\Blk}\bigr)
  +\bigotimes_{\Blk\in X}\bigl(\Pi_{E_1,\Blk}+\varepsilon_\Blk G_{o,\Blk}\bigr)
  \succeq_{\vartheta_0}\bigotimes_{\Blk\in X}\Pi_{F_1,\Blk}
  +\bigotimes_{\Blk\in X}\Pi_{F_2,\Blk},
 \]
 and consequently
 $\bigotimes_\Blk(\Pi_{E_0,\Blk}+\Pi_{E_1,\Blk}+\varepsilon_\Blk G_\Blk)
 \succeq_{\vartheta_0}\bigotimes_\Blk\Pi_{F_1,\Blk}+\bigotimes_\Blk\Pi_{F_2,\Blk}$.
\item Let $C\ge C_X$, $\varepsilon_X\ge\prod_\Blk(1+\varepsilon_\Blk)-1$, and
 suppose $Y[\bigotimes_\Blk\Pi_{F_j,\Blk}]\ge(l_0-o(1))Y$ for $j=1$ and $j=2$.
 Then
 \begin{align*}
  \tfrac12\bigl(Y[E^X_0]+Y[E^X_1]\bigr)&\ge(l_0-C\varepsilon_X-o(1))Y,\\
  2^{-k}\textstyle\sum_{\mathbf s}Y[E_{\mathbf s}]
  &\ge(2^{1-k}l_0-C\varepsilon_X-o(1))Y.
 \end{align*}
\item Under the hypotheses of (b), let $m^X_s$ be the actual half marginal of
 $E^X_s$ at $\iota$. Then
 \[
  Y[\operatorname{Diag}_\iota(m^X_0)]+Y[\operatorname{Diag}_\iota(m^X_1)]
  \ge2\bigl([l_0-C\varepsilon_X]_+-o(1)\bigr)Y .
 \]
\end{enumerate}
By \cref{lem:record:restoration} one may take for $l_0$ the restoration bound
of $X$: the same $h_X$, $M_X$, $\chi_X$ serve $j=1$ and $j=2$, because the
half law of the root does not depend on the chamber (also for the positive
roots, \cref{lem:record:root-marginal}(iv)). No lower bound is asserted for an
individual matching embedding, nor for the diagonal of an individual aligned
embedding.
\end{lemma}

\begin{proof}
(a) Both orders are \cref{lem:tr:polarization}(c), applied to the two-layer
five-type data of the blocks of $X$, which satisfy the hypotheses of
\cref{lem:tr:polarization}(b) with the constants $\varepsilon_\Blk$ (proof of
\cref{lem:record:paired-e-bridge}(a)). The first is the compression to the two
aligned corners of the tensor product of the one-block brackets, the sign of
the fold part being reversed on the corner $\mathcal V$; the second adds the
$2^k-2$ nonaligned corner products
$\bigotimes_\Blk(\Pi_{E_{s_\Blk},\Blk}+\varepsilon_\Blk G_{s_\Blk,\Blk})$, which are
tensor products of positive semidefinite diagonal blocks of these brackets.
The order is not obtained by tensoring a one-sided order with an indefinite
factor.

(b) Apply \ref{it:record:E2} at $\vartheta_0$, with $\mu$ outside, to the two
orders of (a). In the first order, expanding
$\bigotimes_\Blk(\Pi_{E_s,\Blk}+\varepsilon_\Blk G_{s,\Blk})$ gives $E^X_s$ and
products that contain at least one merge law; these are positive measures with
uniform one-slot marginals of total mass $\prod_\Blk(1+\varepsilon_\Blk)-1$, so by
\ref{it:record:E1}
\[
 Y[E^X_0]+Y[E^X_1]\ge2(l_0-o(1))Y-2C\Bigl(\prod_\Blk(1+\varepsilon_\Blk)-1\Bigr)Y.
\]
In the second order the merge products have total mass
$\prod_\Blk(2+2\varepsilon_\Blk)-2^k$, so
$\sum_{\mathbf s}Y[E_{\mathbf s}]\ge2(l_0-o(1))Y-2^kC(\prod_\Blk(1+\varepsilon_\Blk)
-1)Y$. Divide by $2$, respectively $2^k$.

(c) By \cref{lem:record:paired-e-bridge}(b), $E^X_s$ is a positive law symmetric
at $\iota$, a probability coupling of two copies of $m^X_s$; so the contraction
\ref{it:record:P1}, applied to the whole product, gives
$Y[E^X_s]\le Y[\operatorname{Diag}_\iota(m^X_s)]$, and no one-sided order is
tensored with an indefinite factor. Adding the two bounds and using the first
bound of (b) gives the claim with $l_0-C\varepsilon_X$ in place of its positive
part; and $Y[\operatorname{Diag}_\iota(m^X_s)]\ge0$ by \ref{it:record:E3} and
\ref{it:record:E2}.
\end{proof}

The factor $2^{1-k}$ in the second bound of (b) is a genuine loss: only the two
aligned corners carry the restoration. The first bound has no such loss, and
its merge correction is $2\varepsilon_X$ against $2^k\varepsilon_X$ for the second.
For two blocks, $l_0-C\varepsilon_X$ exceeds twice $l_0/2-C\varepsilon_X$, which is why
the aligned gain is at least twice the four-embedding gain
(\cref{cor:var-classes}), and why the construction credits only the aligned
matching classes. Part (c) is the diagonal-mean coupon $2[l_0-C\varepsilon_X]_+$
of the matching diagonals of the mean credit (\cref{prop:record:mean-credit}),
for any number of blocks.

\subsection{The saving $g$: product energy at least $gC_XY$}\label{subsec:variance:gain}

A \emph{selection} of the record expansion (\cref{lem:record:record-expansion})
puts on each block of a finite set $X$ one changed record and one parity,
symmetric or skew, with an even number of skew parities; $I_\omega$ is its
term in $Y_\nu-Y_\lambda$ and $\pi_\omega$ its price, the product over $\Blk\in X$
of the prices of its records with their parities. For a record assignment
$\boldsymbol\phi=(\phi_\Blk)_{\Blk\in X}$ of eligible records put
$\Pi_{\boldsymbol\phi}=\bigotimes_\Blk\Pi_{\phi_\Blk,\Blk}$,
$m_{\boldsymbol\phi}=\bigotimes_\Blk m_{\phi_\Blk,\Blk}$ and
$J_{\boldsymbol\phi}=\bigotimes_\Blk J_{\phi_\Blk,\Blk}$.

\begin{lemma}[Variance gain]\label{lem:var-gain}
Let $X$ be a nonempty finite set of middle blocks, with a parity for each
$\Blk\in X$, an even number of them skew. Let $\mathcal R$ be a nonempty finite
set of record assignments $\boldsymbol\phi$ of eligible records, and let
$\omega_{\boldsymbol\phi}$ be the selection with the records $\phi_\Blk$ and the
fixed parities. Suppose that, for every $\Blk$, the price of $\phi_\Blk$ with
the parity of $\Blk$ is the same for all $\boldsymbol\phi\in\mathcal R$. Let
$C\ge C_X$, $l_0\in\R$ and $0\le\kappa<1$, and assume:
\begin{enumerate}[label=(\textup{\Roman*})]
\item[\textup{(Res)}] $|\mathcal R|^{-1}\sum_{\boldsymbol\phi\in\mathcal R}
 Y[\Pi_{\boldsymbol\phi}]\ge(l_0-o(1))Y$;
\item[\textup{(Cen)}] for every $\boldsymbol\phi\in\mathcal R$, $\Pi_{\boldsymbol\phi}$
 is symmetric at $\iota$, and its Markov operator $T_{\boldsymbol\phi}$ on
 $L^2(m_{\boldsymbol\phi})$ satisfies $\|T_{\boldsymbol\phi}-\Jproj\|\le\kappa$.
\end{enumerate}
Put $\bar l=(l_0/C-\kappa)/(1-\kappa)$ and $g=[\bar l-R_{\rm err}]_+$. Then:
\begin{enumerate}[label=(\alph*)]
\item $|\mathcal R|^{-1}\sum_{\boldsymbol\phi}Y[J_{\boldsymbol\phi}]\ge
 (\bar l-o(1))C_XY$;
\item $\sum_{\boldsymbol\phi}|I_{\omega_{\boldsymbol\phi}}|\le(B-g+o(1))
 \sum_{\boldsymbol\phi}\pi_{\omega_{\boldsymbol\phi}}Y$;
\item if $l_0\le C$, then $0\le g\le1\le B$, so the selections
 $\omega_{\boldsymbol\phi}$, $\boldsymbol\phi\in\mathcal R$, form a class of cost
 $B-g$ in the sense of (CP5).
\end{enumerate}
Any upper bound for $R_{\rm err}$ may replace it in $g$.
\end{lemma}

\begin{proof}
(a) Fix $\boldsymbol\phi$. By (Cen), the centred interpolation \ref{it:record:P3}
at $\iota$, with $\mu$ outside and $\kappa_0=\kappa$, gives
$Y[\Pi_{\boldsymbol\phi}]\le(1-\kappa)Y[J_{\boldsymbol\phi}]+
\kappa Y[\operatorname{Diag}_\iota(m_{\boldsymbol\phi})]$, as
$J_{\boldsymbol\phi}=J_\iota(m_{\boldsymbol\phi})$; and \ref{it:record:E1} bounds
the last energy by $(1+o(1))C_XY$. Averaging over $\mathcal R$ and using (Res),
$(1-\kappa)|\mathcal R|^{-1}\sum_{\boldsymbol\phi}Y[J_{\boldsymbol\phi}]\ge
(l_0-\kappa C_X-o(1))Y$. If $\bar l\le0$, (a) holds because
$Y[J_{\boldsymbol\phi}]\ge0$ by \ref{it:record:E2}, the cut operator of
$J_{\boldsymbol\phi}$ being a projection. If $\bar l>0$, then $l_0>0$, so
$l_0\ge l_0C_X/C$ and $l_0-\kappa C_X\ge(1-\kappa)\bar lC_X$.

(b) For $\boldsymbol\phi\in\mathcal R$ and $\Blk\in X$, the change of $\phi_\Blk$
with the parity of $\Blk$ has, relative to $m_{\phi_\Blk,\Blk}$, a cut operator
$D_\Blk$ at $\iota$ that is self-adjoint or skew-adjoint, annihilates the
constants on both sides by \eqref{eq:record:zero-marginal}, and has norm at most
the coefficient $\tilde\pi_\Blk$ of its norm envelope (\cref{def:record:comparison-law});
by hypothesis $\tilde\pi_\Blk$ does not depend on $\boldsymbol\phi$, and
$\pi_{\omega_{\boldsymbol\phi}}=\prod_\Blk\tilde\pi_\Blk\,C_X$. By
\cref{lem:record:variance-interface}(i),(ii), applied with these $D_\Blk$,
\[
 |I_{\omega_{\boldsymbol\phi}}|\le\prod_\Blk\tilde\pi_\Blk\bigl(
 Y^{\Nbr}[\operatorname{Diag}_\iota(m_{\boldsymbol\phi})]-
 Y^{\Nbr}[J_{\boldsymbol\phi}]\bigr)
 \le\prod_\Blk\tilde\pi_\Blk\bigl((B+R_{\rm err}+o(1))C_XY
 -Y[J_{\boldsymbol\phi}]\bigr).
\]
Summing over $\mathcal R$ and using (a) gives
$\sum|I_{\omega_{\boldsymbol\phi}}|\le(B-(\bar l-R_{\rm err})+o(1))
\sum\pi_{\omega_{\boldsymbol\phi}}Y$. If $\bar l-R_{\rm err}\le0$ use instead
$|I_{\omega_{\boldsymbol\phi}}|\le\prod_\Blk\tilde\pi_\Blk\,
Y^{\Nbr}[\operatorname{Diag}_\iota(m_{\boldsymbol\phi})]\le(B+o(1))
\pi_{\omega_{\boldsymbol\phi}}Y$, valid because $Y^{\Nbr}[J_{\boldsymbol\phi}]\ge0$
by \ref{it:record:E2}. A larger $R_{\rm err}$ only lowers $g$.

(c) If $l_0\le C$ then $\bar l\le1$, so $g\le1$; and
$B=e^{U_0}+R_{\rm err}\ge1$.
\end{proof}

The lemma uses the parities only through the evenness of the number of skew
factors. It applies to two symmetric parities as well as to two skew ones, and
selections with different parities, records or blocks are different monomials
of the record expansion.

\begin{corollary}[The eligible variance classes]\label{cor:var-classes}
Let $X$ be a set of $k\ge1$ middle blocks with fixed parities as in
\cref{lem:var-gain}, and let $C\ge C_X$. Let $h_X$, $M_X$, $\chi_X$ be as in
\cref{lem:record:restoration}, with $0<l=1-(C-1)h_X\le C$, and let $l_0$ be the
largest of $0$ and the three expressions of that lemma. Let
$\kappa_F\ge\max_{\Blk\in X}\kappa_F(\Blk)$ and $\kappa_E\ge\max_{\Blk\in X}
\kappa_E(\Blk)$ be smaller than one, and $\varepsilon_X\ge\prod_{\Blk\in X}(1+
\varepsilon_\Blk)-1$. Assume that the two matching records of every block of $X$
carry the same prices, as they do in the comparison law
(\cref{subsec:variance:eligible}). Then \cref{lem:var-gain} applies to
\begin{enumerate}[label=(\roman*)]
\item the aligned fold assignment $(F_j)_{\Blk\in X}$, $j\in\{1,2\}$, with
 $l_0$ and $\kappa_F$, so that $\bar l_F=(l_0/C-\kappa_F)/(1-\kappa_F)$;
\item the two aligned matching assignments, with $l_0-C\varepsilon_X$ and
 $\kappa_E$, so that $\bar l^{\rm al}_E=(l_0/C-\varepsilon_X-\kappa_E)/(1-\kappa_E)$;
\item the $2^k$ matching embeddings, with $2^{1-k}l_0-C\varepsilon_X$ and
 $\kappa_E$, so that $\bar l_E=(2^{1-k}l_0/C-\varepsilon_X-\kappa_E)/(1-\kappa_E)$.
\end{enumerate}
For $k=1$, (ii) and (iii) coincide; in general
$[\bar l^{\rm al}_E-R_{\rm err}]_+\ge2^{k-1}[\bar l_E-R_{\rm err}]_+$.
\end{corollary}

\begin{proof}
(i) (Res) is \cref{lem:record:restoration}, where the maximum with $0$ is
allowed because $Y[\bigotimes_\Blk\Pi_{F_j,\Blk}]\ge0$; (Cen) is
\cref{lem:record:old-f-marginal}(d). (ii), (iii) (Res) is
\cref{lem:record:paired-e-joint}(b), with the restoration bound for $j=1$ and
$j=2$; (Cen) is \cref{lem:record:paired-e-bridge}(b). The prices are common by
assumption in (ii), (iii) and trivially in (i). Finally
$2^{k-1}\bar l_E=(l_0/C-2^{k-1}(\varepsilon_X+\kappa_E))/(1-\kappa_E)\le
\bar l^{\rm al}_E$, and $R_{\rm err}\ge0$, $2^{k-1}\ge1$ give the last claim.
\end{proof}

\begin{remark}[Repair independence]\label{rem:variance:independence}
The numbers $l_0$, $\kappa_F$, $\kappa_E$ and $\varepsilon_X$, hence the gains,
depend only on the reference, its records, the bracket (through $B$ and
$R_{\rm err}$) and the cap $C$. A repair enters \cref{lem:var-gain} only through
the zero-marginal identity \eqref{eq:record:zero-marginal} and through its
prices. The credits therefore hold for every repair in the sense of
\cref{def:record:eligible}, signed or not, in particular for the repairs of the
final construction.
\end{remark}

Selections that put $F_1$ on one block and $F_2$ on another, or a fold record
on one block and a matching record on another, receive no variance credit:
$F_1$ and $F_2$ are diagonal at different chamber cuts, so no chamber diagonal
is their product, and the restoration needs a common chamber. Their skew
classes are priced by positive envelopes (\cref{thm:envelopes}).

\subsection{Additional positive terms in the exterior energy}
\label{subsec:variance:backgrounds}

The gains of \cref{lem:var-gain} use the selected blocks only. Outside them the
exterior bracket $\Nbrz$ exceeds $\mu$ by positive reserved-diagonal records
(\cref{def:record:frame}, \cref{lem:record:gram}(i)), and the variance form may
use the energy $Y^{\Nbrz}[J_X]$ with $\Nbrz$ outside in place of $Y[J_X]$, at the
cost of one comparison of $\Nbr$ with $\Nbrz$
(\cref{lem:record:variance-interface}(iii)). Among these records are the two
matching diagonals of $E_{\rm add}$ of every middle block, which are diagonal
copies of the actual half marginals of the two matching records. A selected
aligned matching product therefore gains, from every set of other blocks
carrying $E_{\rm add}$, the energy of its product of actual half marginals
tensored with the diagonals of these blocks, and a joint restoration of the
selected and the other blocks bounds this energy below. This subsection
bounds these \emph{positive background} terms: on one or two other blocks of any
pieces, with a lower bound for the weighted number of distinct choices of these
blocks, and on three or four blocks of a pool of late pieces, where uniform
data suffice. The bracket has one fold copy, which is a summand of $\Nbr$ and
not of $\Nbrz$ (\cref{lem:record:gram}); so the exterior contains no fold
diagonal, and a selected fold variance has no positive background term of this
kind.

\paragraph{Piece data.}
The bounds are stated for \emph{piece data}: the middle blocks are grouped into
the pieces $[X_i,Y_i)$, every block of the piece $i$ has least prime at least
$X_i$, and numbers $h_i\ge0$, $M_i\ge1$, $\chi_i\ge0$, $\varepsilon_i\ge0$,
$k_{F,i},k_{E,i}\in[0,1)$, $C_i\ge1$ and $\Sigma_i\ge0$ satisfy, for every block
$\Blk$ of the piece $i$,
\begin{gather*}
 h_\Blk\le h_i,\qquad M_\Blk\le M_i,\qquad
 \chi^2(\nu_{\mathrm L,\Blk}\|\lambda_{F,\Blk})\le\chi_i,\qquad
 \varepsilon_\Blk\le\varepsilon_i,\\
 \kappa_F(\Blk)\le k_{F,i},\qquad\kappa_E(\Blk)\le k_{E,i},\qquad
 q_\Blk^{2a}\le C_i,
\end{gather*}
and $\Sigma_i\le\sum_{\Blk\subseteq[X_i,Y_i)}b_\Blk^{-1}$. For a finite multiset $m$
of pieces, with multiplicities $m_i$ and size $|m|$, put
\begin{gather*}
 h_m=\prod_{i\in m}(1+h_i)-1,\quad M_m=\prod_{i\in m}M_i,\quad
 \chi_m=\prod_{i\in m}(1+\chi_i)-1,\quad\eta_m=\prod_{i\in m}(1+\varepsilon_i)-1,\\
 C_m=\prod_{i\in m}C_i,\qquad k_{F,m}=\max_{i\in m}k_{F,i},\qquad
 k_{E,m}=\max_{i\in m}k_{E,i},\qquad\ell_m=\mathcal F(C_m,h_m,M_m,\chi_m),
\end{gather*}
the products with multiplicity, $C_\varnothing=1$, and $m+r$ the union of two
multisets. A set $X$ of distinct middle blocks has the \emph{pattern} $m$ if
$m_i$ of its blocks lie in the piece $i$. Then the products over $X$ of
$1+h_\Blk$, of $M_\Blk$, of $1+\chi^2(\nu_{\mathrm L,\Blk}\|\lambda_{F,\Blk})$ and of
$1+\varepsilon_\Blk$ are at most $1+h_m$, $M_m$, $1+\chi_m$ and $1+\eta_m$, and
$C_X\le C_m$; so \cref{lem:record:restoration} gives
$Y[\bigotimes_{\Blk\in X}\Pi_{F_j,\Blk}]\ge(\ell_m-o(1))Y$ for $j=1,2$, and
\cref{cor:var-classes} applies with $C=C_m$, $\kappa_F=k_{F,m}$,
$\kappa_E=k_{E,m}$ and $\varepsilon_X=\eta_m$ (if $1-(C_m-1)h_m\le0$, then
$\ell_m=0$ and \cref{lem:var-gain} applies directly with $l_0=0$). Its numbers
(i), (ii) are the \emph{margins without background blocks}
(called selected-only margins in the data)
\begin{equation}\label{eq:variance:margins}
 b_F(m)=\frac{\ell_m/C_m-k_{F,m}}{1-k_{F,m}},\qquad
 b_E(m)=\frac{\ell_m/C_m-\eta_m-k_{E,m}}{1-k_{E,m}} .
\end{equation}
The coefficients of the two matching diagonals of $E_{\rm add}$, divided by the
normalizer $\widetilde\mu_\Blk(1)=1+h_\Blk$, are $v_{{\rm D}}r_\Blk^2/(1+h_\Blk)$
(\cref{def:record:frame}(b)); in the piece $i$ they are at least
$c_{E,i}b_\Blk^{-1}$, where
\[
 c_{E,i}=\frac{v_{{\rm D},i}}{1+h_i},
\]
with $v_{{\rm D},i}$ the completion weight of the piece.

\begin{lemma}[Distinct choices]\label{lem:distinct-choice}
Let $c=(c_i)$ be nonnegative numbers indexed by the pieces, $m$ a multiset of
pieces, $X$ a set of distinct middle blocks of pattern $m$, and write $i(\Blk)$
for the piece of a block $\Blk$.
\begin{enumerate}[label=(\alph*)]
\item For a multiset $r$ of pieces, the sum of
 $\prod_{\Blk\in Z}c_{i(\Blk)}b_\Blk^{-1}$ over the sets $Z$ of middle blocks of
 pattern $r$ disjoint from $X$ is at least the choice weight $W(m,r;c)$ (not the wrapping term $W$),
 \begin{equation}\label{eq:variance:choice-weight}
  W(m,r;c)=\prod_i\frac{c_i^{r_i}}{r_i!}\prod_{j=0}^{r_i-1}
  \Bigl(\Sigma_i-\frac{m_i+j}{X_i}\Bigr)_+ .
 \end{equation}
\item Let $\mathcal P$ be a set of pieces and put
 $\mathfrak m(c)=\sum_{i\in\mathcal P}c_i\Sigma_i$,
 $\mathfrak u(c)=\max_{i\in\mathcal P}c_i/X_i$ and
 $\mathfrak e_m(c)=\sum_{i\in\mathcal P}m_ic_i/X_i$. For $k\ge1$ the same sum over
 the sets $Z$ of $k$ distinct blocks of pieces of $\mathcal P$, disjoint from $X$,
 is at least
 \[
  \mathfrak w_k(c;m)=\frac1{k!}\prod_{j=0}^{k-1}
  \bigl[\mathfrak m(c)-\mathfrak e_m(c)-j\,\mathfrak u(c)\bigr]_+ .
 \]
\end{enumerate}
\end{lemma}

\begin{proof}
Every block of the piece $i$ has $b_\Blk^{-1}\le1/X_i$. (a) The choices in
different pieces are independent, so it suffices to treat one piece $i$, and to
count ordered $r_i$-tuples of distinct blocks, then divide by $r_i!$. The
blocks of the piece outside $X$ have $\sum b_\Blk^{-1}\ge\Sigma_i-m_i/X_i$, since
$X$ removes $m_i$ of them; after $j$ more choices the remaining blocks have
$\sum b_\Blk^{-1}\ge\Sigma_i-(m_i+j)/X_i$. Summing successively from the last
choice, each partial sum is at least the positive part of its bound, whatever
the earlier choices; multiply by $c_i^{r_i}$. (b) The same argument on the
union of the pieces of $\mathcal P$: the blocks of $\mathcal P$ outside $X$ have
$\sum c_{i(\Blk)}b_\Blk^{-1}\ge\mathfrak m(c)-\mathfrak e_m(c)$, and each later
choice removes a block of weight $c_{i(\Blk)}b_\Blk^{-1}\le\mathfrak u(c)$;
divide by $k!$.
\end{proof}

\paragraph{Pooling the pieces from $2^{31}$.}
Let $\mathcal P$ be the set of middle pieces with $X_i\ge2^{31}$, and give it the
\emph{pool data} $h_{\mathcal P}=\max_{i\in\mathcal P}h_i$ and likewise
$M_{\mathcal P}$, $\chi_{\mathcal P}$, $\varepsilon_{\mathcal P}$, and the cap
$C_{\mathcal P}=\max_{i\in\mathcal P}C_i$. For a multiset $m$ and $k\ge1$ let
$m+\mathcal P^k$ denote $m$ with $k$ additional entries carrying the pool data, so
that $h_{m+\mathcal P^k}$, $M_{m+\mathcal P^k}$, $\chi_{m+\mathcal P^k}$,
$\eta_{m+\mathcal P^k}$, $C_{m+\mathcal P^k}=C_mC_{\mathcal P}^k$ and
$\ell_{m+\mathcal P^k}$ are defined as above, and write $c^{\mathcal P}$ for a
coefficient vector $c$ restricted to $\mathcal P$ (zero off $\mathcal P$).

\begin{lemma}[Terms pooled over the pieces from $2^{31}$]\label{lem:pool-coupons}
If $X$ has pattern $m$ and $Z$ is a set of $k$ distinct blocks of pieces of
$\mathcal P$ disjoint from $X$, then the four products over $X\cup Z$ of the
piece-data paragraph are at most those of $m+\mathcal P^k$, and
$C_{X\cup Z}\le C_mC_{\mathcal P}^k$; so
$Y[\bigotimes_{\Blk\in X\cup Z}\Pi_{F_j,\Blk}]\ge(\ell_{m+\mathcal P^k}-o(1))Y$ for
$j=1,2$, and \cref{lem:record:paired-e-joint}(b), (c) hold for $X\cup Z$ with
$l_0=\ell_{m+\mathcal P^k}$, $C=C_mC_{\mathcal P}^k$ and
$\varepsilon_{X\cup Z}\le\eta_{m+\mathcal P^k}$. Moreover:
\begin{enumerate}[label=(\alph*)]
\item \emph{Variance.} For $|m|\le2$ and an upper bound $R_*\ge R_{\rm err}$ put
 \[
  P_E(m)=\sum_{k\in\{3,4\}}\mathfrak w_k(c_E^{\mathcal P};m)\,
  \frac{[\ell_{m+\mathcal P^k}/C_m-C_{\mathcal P}^k(\eta_{m+\mathcal P^k}+k_{E,m})]_+}
  {1-k_{E,m}}
 \]
 if $b_E(m)\ge R_*$, and $P_E(m)=0$ otherwise. This term may be added to the gain
 of the aligned matching variance of \cref{prop:background-variance}, whose proof
 shows it.
\item \emph{Mean.} Let $k\in\{5,6\}$ and let $Z$ range over the sets of $k$
 distinct blocks of pieces of $\mathcal P$. For coefficient vectors $d,d'$ with
 $0\le d_i\le\rho\,d'_i$ on $\mathcal P$, for some $\rho\in[0,1]$,
 \[
  \sum_Z\prod_{\Blk\in Z}d_{i(\Blk)}b_\Blk^{-1}\ge\mathfrak w_k(d^{\mathcal P};\varnothing),
  \qquad
  \sum_Z\Bigl[\prod_{\Blk\in Z}d'_{i(\Blk)}b_\Blk^{-1}-\prod_{\Blk\in Z}
  d_{i(\Blk)}b_\Blk^{-1}\Bigr]\ge(1-\rho^k)\,\mathfrak w_k(d'^{\mathcal P};\varnothing).
 \]
 With the restorations of the first assertion, these give the late-pool
 coupons of the mean credit (\cref{prop:record:mean-credit}),
 \begin{align*}
  m^{\rm pool}_F&=\sum_{k\in\{5,6\}}\ell_{\mathcal P^k}\Bigl[2\mathfrak w_k(c_F^{\mathcal P};\varnothing)
  +\mathfrak w_k\bigl((k_{\rm L}/H)^{\mathcal P};\varnothing\bigr)
  +\mathfrak w_k\bigl((k_{\rm R}/H)^{\mathcal P};\varnothing\bigr)\Bigr],\\
  m^{\rm pool}_Q&=\sum_{k\in\{5,6\}}2(1-\rho^k)\,
  \mathfrak w_k\bigl((c_E+q)^{\mathcal P};\varnothing\bigr)
  \bigl[\ell_{\mathcal P^k}-C_{\mathcal P}^k\eta_{\mathcal P^k}\bigr]_+,
 \end{align*}
 with $c_{F,i}=c_{E,i}/H_i$, the density factors $H_i$ of
 \eqref{eq:variance:pieces}, the side coefficients $k_{s,i}$ of
 \cref{prop:signed-sides}, the coefficients $q_i$ of the chosen matching
 diagonals (\cref{subsec:matching-envelopes}) and
 $\rho=\max_{i\in\mathcal P}c_{E,i}/(c_{E,i}+q_i)$. They occupy the degrees five
 and six, disjoint from the degrees at most four of the other coupons.
\end{enumerate}
\end{lemma}

\begin{proof}
Every block of a piece of $\mathcal P$ has data at most the pool data and
$q_\Blk^{2a}\le C_i\le C_{\mathcal P}$, so the products are bounded as stated;
the pool data are therefore valid hypotheses of \cref{lem:record:restoration}
and of \cref{lem:record:paired-e-joint}(b), (c), whose conclusions follow.
(a) The bound that justifies the term is part of the proof of
\cref{prop:background-variance}. If $b_E(m)\ge R_*$ then
$b_E(m)-R_*\ge0$, so after the positive part of the remaining
gain is taken, this difference is exactly the term the
certificate adds; the complementary term $[R_*-b_E(m)]_+$ is
then zero. (b) The first bound is
\cref{lem:distinct-choice}(b) with $m$ empty. For the second,
$\prod_{\Blk\in Z}d_{i(\Blk)}\le\rho^k\prod_{\Blk\in Z}d'_{i(\Blk)}$ pointwise, so the
bracket is at least $(1-\rho^k)\prod_{\Blk\in Z}d'_{i(\Blk)}b_\Blk^{-1}$; apply the
first bound to $d'$. In the coupons the fold diagonals of $k$ pool blocks have,
by the first assertion and \cref{lem:record:companion-energy}(a) (the energy bound for the
one-layer companions $D^j_\Blk$ of the fold records, proved in
\cref{sec:matching-mean}), applied with the set $Z$, energy at least
$\ell_{\mathcal P^k}$ divided by the product of the $H_i$, and the two aligned
matching diagonals have energies
adding up to at least $2[\ell_{\mathcal P^k}-C_{\mathcal P}^k\eta_{\mathcal P^k}]_+$
(\cref{lem:record:paired-e-joint}(c)). The coefficient per block of the polarized
fold branch is at least $c_{E,i}b_\Blk^{-1}$ and that of a chosen side diagonal
$k_{s,i}b_\Blk^{-1}$, which with the factors $1/H_i$ of the energies give the
vectors $c_F$ and $k_s/H$; the nonempty-$Q$ branch (the terms in the expansion of
$\prod_{\Blk\in Z}(c_{E,i(\Blk)}+q_{i(\Blk)})b_\Blk^{-1}$ in which the
$q$-factor is chosen for at least one block) has the coefficient
$\prod(c_E+q)b^{-1}-\prod c_Eb^{-1}$, to which the second bound applies with
$d=c_E$, $d'=c_E+q$. The assembly of the branches is that of
\cref{prop:record:mean-credit}.
\end{proof}

\begin{proposition}[Positive backgrounds in a selected variance]
\label{prop:background-variance}
Assume the piece data above. For multisets $m$, $r$ of pieces put
\[
 p_E(m,r)=\frac{[\ell_{m+r}/C_m-C_r(\eta_{m+r}+k_{E,m})]_+}{1-k_{E,m}},
\]
and for $|m|\in\{1,2\}$ and an upper bound $R_*\ge R_{\rm err}$
\begin{equation}\label{eq:variance:background-gain}
 g_E(m)=\Bigl[b_E(m)+\sum_{1\le|r|\le2}W(m,r;c_E)\,p_E(m,r)+P_E(m)-R_*\Bigr]_+,
 \qquad g_F(m)=\bigl[b_F(m)-R_*\bigr]_+,
\end{equation}
with $P_E(m)$ of \cref{lem:pool-coupons}(a). Let $X$ be a set of distinct middle
blocks of pattern $m$, with parities, an even number of them skew.
\begin{enumerate}[label=(\alph*)]
\item With $J_s$ the product of the actual half marginals of the aligned
 embedding $E^X_s$,
 \[
  \tfrac12\bigl(Y^{\Nbrz}[J_0]+Y^{\Nbrz}[J_1]\bigr)\ge
  \Bigl(C_m\Bigl[b_E(m)+\sum_{1\le|r|\le2}W(m,r;c_E)p_E(m,r)+P_E(m)\Bigr]-o(1)\Bigr)Y .
 \]
 If $g_E(m)\le1$, the class of the two aligned matching assignments
 $(E_0)_{\Blk\in X}$, $(E_1)_{\Blk\in X}$ with these parities has the class cost
 $B-g_E(m)$ in the sense of (CP5); in $g_E(m)$ the error $R_*$ of comparing
 $\Nbr$ with $\Nbrz$ outside $X$ is subtracted once, for all the
 kept terms together.
\item For $j\in\{1,2\}$ the selection of $F_j$ on $X$ with these parities has the
 class cost $B-g_F(m)$.
\end{enumerate}
\end{proposition}

\begin{proof}
(a) Fix $s\in\{0,1\}$ and write $D_{Z,s}=\bigotimes_{\Blk\in Z}D_{E_s,\Blk}$ for a
set $Z$ of middle blocks disjoint from $X$, with the selected diagonals
$D_{E_s,\Blk}=\operatorname{Diag}_\iota(m_{E_s,\Blk})$ of
\cref{def:record:eligible}.

\emph{Exterior expansion.} On a middle block $\Blk\notin X$, $(\Nbrz-\mu)_\Blk$ is a
positive reserved-diagonal measure (\cref{lem:record:gram}(i)), and it contains
the matching diagonal of $E_{\rm add}$ of type $s$: the normalized measure
$(\Nbrz-\mu)_\Blk-v_{\rm D}r_\Blk^2(1+h_\Blk)^{-1}D_{E_s,\Blk}$ is again positive
and reserved diagonal, because the half law of $\Pi^{(1)}_{E_s}$ is $m_{E_s,\Blk}$
(\cref{def:record:frame}(b) and the paragraph after it); outside the middle
$\Nbr_{0,\Blk}=\mu_\Blk$. Expand $\Nbr_{0,-X}=\bigotimes_{\Blk\notin X}(\mu_\Blk+(\Nbrz-\mu)_\Blk)$
over the sets $Z$ of middle blocks outside $X$. The term of $Z$ is the energy, with
$\mu$ outside $X\cup Z$, of $J_s\otimes(\Nbrz-\mu)_Z$, a tensor product of
measures $\succeq_\iota0$ ($J_s$ has a projection as cut operator, the others are
positive and reserved diagonal, \ref{it:record:E3}); by \ref{it:record:E2} it is
nonnegative, and, since $v_{\rm D}r_\Blk^2(1+h_\Blk)^{-1}\ge c_{E,i(\Blk)}b_\Blk^{-1}$
($r_\Blk^2=b_\Blk^{-1}$, $h_\Blk\le h_i$) and the difference of the two tensor
products is a sum of tensor products of measures $\succeq_\iota0$, it is at least
$\prod_{\Blk\in Z}c_{E,i(\Blk)}b_\Blk^{-1}\cdot Y[J_s\otimes D_{Z,s}]$. We keep the
term $Z=\varnothing$, the terms with $1\le|Z|\le2$, and, if $b_E(m)\ge R_*$, the
terms with $Z$ a set of three or four blocks of the pool; they are distinct terms
of the expansion, and the others are dropped. The middle range has finitely
many blocks, so the $o(1)$ below add up to $o(1)$.

\emph{The selected term.} By \cref{lem:record:paired-e-bridge}(b), $E^X_s$ is
symmetric at $\iota$ with centred norm at most $k=k_{E,m}$, so \ref{it:record:P3}
gives $(1-k)Y[J_s]\ge Y[E^X_s]-k(1+o(1))C_XY$. Averaging over $s$, using
\cref{lem:record:paired-e-joint}(b) with $l_0=\ell_m$, $C=C_m$, $\varepsilon_X\le\eta_m$,
and $C_X\le C_m$, gives $\frac12\sum_sY[J_s]\ge(C_mb_E(m)-o(1))Y$.

\emph{A background term.} Let $Z$ have pattern $r$. By \ref{it:record:P3b} with
$\pi=E^X_s$ ($\kappa_0=k$) and $\pi'=E^Z_s$, whose half marginal at $\iota$ is
$\bigotimes_{\Blk\in Z}m_{E_s,\Blk}$,
\[
 Y[E^{X\cup Z}_s]\le(1-k)\,Y[J_s\otimes D_{Z,s}]+k\,Y[\operatorname{Diag}_\iota(
 m^X_s\otimes m^Z_s)]\le(1-k)\,Y[J_s\otimes D_{Z,s}]+k(1+o(1))C_{X\cup Z}Y,
\]
where $E^{X\cup Z}_s=E^X_s\otimes E^Z_s$ is the aligned embedding of $X\cup Z$. The
set $X\cup Z$ has pattern $m+r$, so \cref{lem:record:paired-e-joint}(b) applies to it
with $l_0=\ell_{m+r}$, $C=C_mC_r\ge C_{X\cup Z}$ and $\varepsilon_{X\cup Z}\le\eta_{m+r}$.
Averaging over $s$, and using $Y[J_s\otimes D_{Z,s}]\ge0$ (\ref{it:record:E2}),
\[
 \tfrac12\sum_sY[J_s\otimes D_{Z,s}]\ge(C_mp_E(m,r)-o(1))Y .
\]
The coefficient $\prod_{\Blk\in Z}c_{E,i(\Blk)}b_\Blk^{-1}$ is the same for $s=0$ and
$s=1$, and its sum over the sets $Z$ of pattern $r$ disjoint from $X$ is at least
$W(m,r;c_E)$ by \cref{lem:distinct-choice}(a). For a set $Z$ of $k'\in\{3,4\}$
blocks of the pool, the same argument with the first assertion of
\cref{lem:pool-coupons} in place of the piece data of $X\cup Z$ gives the bound
with
$[\ell_{m+\mathcal P^{k'}}/C_m-C_{\mathcal P}^{k'}(\eta_{m+\mathcal P^{k'}}+k)]_+/(1-k)$,
and the coefficients sum to at least $\mathfrak w_{k'}(c_E^{\mathcal P};m)$ by
\cref{lem:distinct-choice}(b). Adding the kept terms gives the displayed bound.

\emph{The class cost.} For $s\in\{0,1\}$ let $\omega_s$ be the selection of
$E^X_s$ with the given parities, and $\gamma=\prod_\Blk\tilde\pi_\Blk$ the common
coefficient of its norm envelopes (the two matching records of a block carry the
same prices), so that $\pi_{\omega_s}=\gamma C_X$. By
\cref{lem:record:variance-interface}(i), (ii) and the last bound of
\ref{it:record:E1},
$|I_{\omega_s}|\le\gamma\bigl((B+o(1))C_XY-Y^{\Nbr}[J_s]\bigr)$ and
$Y^{\Nbr}[J_s]\ge Y^{\Nbrz}[J_s]-(R_{\rm err}+o(1))C_XY$; this is the comparison
of $\Nbr$ with $\Nbrz$, made once for all the kept terms together. Let $\tilde b$ be
the bracket in the displayed bound. If $\tilde b>R_*$, then $\tilde b>0$ and
$C_m\tilde b\ge C_X\tilde b$, so summing over $s$,
\[
 |I_{\omega_0}|+|I_{\omega_1}|\le\bigl(B-(\tilde b-R_{\rm err})+o(1)\bigr)
 (\pi_{\omega_0}+\pi_{\omega_1})Y\le\bigl(B-g_E(m)+o(1)\bigr)(\pi_{\omega_0}
 +\pi_{\omega_1})Y .
\]
If $\tilde b\le R_*$, then $g_E(m)=0$, and $|I_{\omega_s}|\le\gamma Y^{\Nbr}[
\operatorname{Diag}_\iota(m^X_s)]\le(B+o(1))\pi_{\omega_s}Y$ as in the proof of
\cref{lem:var-gain}(b). Finally $0\le g_E(m)\le1\le B$.

(b) The exterior contains no fold diagonal, and (b) is \cref{cor:var-classes}(i)
with \cref{lem:var-gain}(c), $C=C_m$ and $\bar l_F=b_F(m)$.
\end{proof}

The background terms lower one selected variance; they are not separate
selected records, and they carry no price of their own. All of them are terms
of the one expansion of $\Nbrz$ outside the selected blocks, so the exterior error
$R_{\rm err}$ is paid once, and the positive part is taken only after it is
subtracted. On the two aligned assignments the background blocks carry the matching
diagonal of the same type as the selected records, and the restoration of the
union keeps both fold products of the aligned bridge; no assertion is made for a
crossed matching assignment. Pooling is restricted to the pieces
from $2^{31}$ because a uniform bound using the maxima over all
pieces would lose the term; over those later pieces the maxima are the pool data with which
the pooled term $P_E(m)$ of \eqref{eq:variance:background-gain} is
formed, and that term is included at the certified piece values of
\cref{cert:variance:fresh} when $b_E(m)\ge R_*$.

\subsection{Assembly: the total saving $D_{\rm var}$}\label{subsec:variance:final}

The final construction uses the credits of degrees one to four. They are
evaluated once per middle piece and per multiset of pieces, at the piece
starts.

\paragraph{Piece values.}
Let $[X_i,Y_i)$, $i=1,\dots,33$, be the middle pieces
(\cref{def:record:exact-blocks}, \cref{tab:record:middle-pieces}). Every block
of the piece $i$ has least prime at least $X_i$ and at most
$\lfloor L/\log X_i\rfloor\le n_i$ primes (\cref{lem:record:local-domain}), where
$L=570$ and $n_i=\lfloor L/\log_-X_i\rfloor$ with a rational lower bound
$\log_-X_i\le\log X_i$; and every coefficient below is a nondecreasing function
of $(q,r)$ (\cref{lem:majorant}), evaluated at $q=1/X_i$ and
$r=1/\lfloor\sqrt{X_i}\rfloor\ge X_i^{-1/2}$, except that on the seven pieces with
positive roots $R_\ulaw$ and $d_\ulaw$ are the endpoint bounds of
\cref{cert:positive-root-caps}, valid at every larger prime. The \emph{piece values}
are upper bounds
\begin{equation}\label{eq:variance:pieces}
\begin{gathered}
 h_i\ge h(X_i),\qquad M_i\ge\bar R(X_i)^{n_i},\qquad
 \chi_i\ge\bigl(1+\bar\chi(X_i)\bigr)^{n_i}-1\qquad(i\le26),\\
 M_i\ge R_\ulaw(X_i)^{n_i},\qquad
 \chi_i\ge\bigl(1+d_\ulaw(X_i)^2\bigr)^{n_i}-1\qquad(27\le i\le33),\\
 k_{F,i}\ge\kappa_F\bigl(1/X_i,1/\lfloor\sqrt{X_i}\rfloor\bigr),\qquad
 k_{E,i}=\frac3{X_i-3},\qquad
 \varepsilon_i\ge\frac{f_2+h_2}{\lfloor\sqrt{X_i}\rfloor},\qquad
 H_i=\Bigl(1+\frac2{X_i-4}\Bigr)^{n_i},
\end{gathered}
\end{equation}
where $h(X_i)$ is \eqref{eq:record:nonroot-mass} with the completion weights of
the piece (\cref{sec:middle-frame}), $\bar R$ and $\bar\chi$ are those of
\cref{lem:record:root-marginal}(iii), $R_\ulaw$ and $d_\ulaw$ those of
\cref{lem:record:root-marginal}(iv) on the seven pieces with positive roots,
$\kappa_F$ is the centred norm of \cref{lem:record:old-f-marginal}(c),
$f_2,h_2$ are the two-layer row coefficients of the piece, and $H_i$ is the
density factor of the actual fold marginal used by the mean savings
(\cref{lem:record:companion-energy}). By \cref{lem:record:root-marginal},
\cref{lem:record:old-f-marginal}, \cref{lem:record:paired-e-bridge}(b) and
monotonicity, every block $\Blk$ of the piece $i$ has $h_\Blk\le h_i$,
$M_\Blk\le M_i$, $\chi^2(\nu_{\mathrm L,\Blk}\|\lambda_{F,\Blk})\le\chi_i$,
$\kappa_F(\Blk)\le k_{F,i}$, $\kappa_E(\Blk)\le k_{E,i}$ and
$\varepsilon_\Blk\le\varepsilon_i$. Let $\Sigma_i$ be the certified lower bound of
\cref{lem:record:lower-start-sums} for $\sum_{\Blk\subseteq[X_i,Y_i)}b_\Blk^{-1}$.
The piece values do not depend on $a$; they are listed in
\cref{tab:data:piece-values}.

\paragraph{The cap.}
Let $a_0=\PushAnchor$ be the evaluation exponent at which the
middle prices are stored in \cref{tab:data:prices},
$t=e^{2(a-a_0)L}$ the repricing factor, and $C_i=e^{2aL_i}$ the piece caps of
\cref{lem:piece-caps} at the anchor $a$, with the authenticated log caps
$L_i\le L$ (the first is $L_1=9135105/16384$); every middle block of the piece $i$
has $q_\Blk^{2a}\le C_i$, and a set of blocks of pattern $m$ has $C_X\le C_m$.
Together with the piece values, the caps and the numbers $\Sigma_i$ are piece data
in the sense of \cref{subsec:variance:backgrounds}. Let $B_*\ge B$ and
$R_*\ge R_{\rm err}$ be the certified upper bounds of
\cref{def:closing-quantities}(2) at the exponent $a$
(\cref{cor:first-stage-repricing,cor:two-cut-series}).

\paragraph{Fresh restorations and gains.}
With these data, $\ell_m=\mathcal F(C_m,h_m,M_m,\chi_m)$ is formed for every
multiset $m$ of pieces of size at most four, and on the pool of size at most six,
the selected-only margins $b_F(m)$, $b_E(m)$ are those of
\eqref{eq:variance:margins}, and the gains are
\begin{equation}\label{eq:variance:fresh}
 g_F(m),\ g_E(m)\ \ \text{of \eqref{eq:variance:background-gain}}\ \ (|m|\le2),
 \qquad
 g^{\rm sel}_T(m)=\bigl[b_T(m)-R_*\bigr]_+\ \ (T\in\{F,E\},\ |m|\le4).
\end{equation}
Thus $g_F=g^{\rm sel}_F$, and $g_E(m)$ adds to $g^{\rm sel}_E(m)$ the positive
backgrounds of degrees one and two and the late-pool terms of degrees three and
four. The restorations depend only on the reference and the caps
(\cref{rem:variance:independence}); they are also the restorations of the mean
coupons (\cref{lem:record:companion-energy,prop:record:mean-credit}).

\paragraph{Prices.}
For the piece $i$ let $S_{F,i}$, $S_{E,i}$ (symmetric), $A_{F,i}$, $A_{E,i}$ (full
skew) and $A^{\rm arm}_{F,i}$, $A^{\rm arm}_{E,i}$ (raw arm) be the prices of
\cref{def:closing-quantities}(3) (\cref{tab:data:prices}): rounded upper bounds,
at the exponent $a_0$, for the sums over the blocks of the piece of the prices of
the two fold and of the two matching records, and of their raw-arm auxiliary
prices, given by \cref{lem:middle-comparison-prices}; with the degree charges
$\bar I_{d,i}$ of the piece,
\begin{gather*}
 S_{F,i}=2w_{{\rm D},i}s_F(X_i)\bar I_{3,i},\qquad
 A_{F,i}=2w_{{\rm D},i}\ell_\tlaw(X_i)\bar I_{2,i}+2w_{{\rm D},i}a_F(X_i)\bar I_{3,i},\qquad
 A^{\rm arm}_{F,i}=2w_{{\rm D},i}\bar I_{2,i},\\
 S_{E,i}=2w_{{\rm S},i}s_E(X_i)\bar I_{2,i},\qquad
 A_{E,i}=2w_{{\rm S},i}a_E(X_i)\bar I_{2,i},\qquad
 A^{\rm arm}_{E,i}=2w_{{\rm S},i}\bar I_{2,i},
\end{gather*}
each rounded up as in \cref{def:closing-quantities}(3). Every middle block has
$\log q_\Blk\le L$, so at the exponent
$a$ these prices are at most $t$ times their values
(\cref{lem:record:repricing}). By \cref{lem:record:record-expansion}(iii) each
piece total is distributed over the blocks of the piece in proportion to their
prices, the two records of one kind in a block receiving the same price
(\cref{subsec:variance:eligible}); these per-block upper prices are the prices of
the complete polynomial.

\paragraph{The credit.}
Put, with the double sums over ordered pairs of pieces and $g(i,j)$ the gain of
the multiset $\{i,j\}$,
\begin{equation}\label{eq:variance:dvar}
\begin{split}
 D_{\rm var}={}&t\sum_i\bigl[g_F(i)S_{F,i}+g_E(i)S_{E,i}\bigr]
 +\frac{t^2}4\sum_{i,j}\Bigl[g_F(i,j)\bigl(S_{F,i}S_{F,j}
 +\kappa_{FF}A^{\rm arm}_{F,i}A^{\rm arm}_{F,j}\bigr)\\
 &+g_E(i,j)S_{E,i}S_{E,j}
 +\kappa^{\rm al}_b\,g^{\rm sel}_E(i,j)A^{\rm arm}_{E,i}A^{\rm arm}_{E,j}\Bigr]\\
 &+\sum_{T\in\{F,E\}}\sum_{3\le|m|\le4}\frac{t^{|m|}g^{\rm sel}_T(m)}
 {2^{|m|-1}\prod_im_i!}
 \Bigl[\sum_\sigma\prod_{l=1}^{|m|}P^{\sigma_l}_{T,i_l}
 +\one_{|m|=4}\,\kappa_{4,T}\prod_{l=1}^4A^{\rm arm}_{T,i_l}\Bigr],
\end{split}
\end{equation}
where $i_1\le\dots\le i_{|m|}$ are the pieces of the multiset $m$, $\sigma$ runs
over the words in $\{{\rm s},{\rm a}\}^{|m|}$ with an even number of letters ${\rm a}$,
not all equal to ${\rm a}$, $P^{\rm s}_{T,i}=S_{T,i}$ and $P^{\rm a}_{T,i}=A_{T,i}$;
and where $\kappa_{FF}$ is the cap of the envelope FF of the fold pairs,
$\kappa^{\rm al}_b$ the invariant cap of the envelope EE aligned of the aligned
matching pairs, $\kappa_{4,F}$ the cap of the envelope F4E0 of the class $(4,0)$
and $\kappa_{4,E}$ the invariant cap of the aligned group of the envelope E4 of
the class $(0,4)$ (\cref{thm:envelopes,lem:class-credit}; \cref{tab:data:envelopes}). The factors
$\frac14$ and $2^{1-|m|}$ account for the two records of one kind of a block,
which carry equal prices, of which the variance classes retain the two aligned
assignments.

\begin{corollary}[The variance credits of the final construction]
\label{cor:variance:final}
Assume the gates of \cref{cert:variance:fresh}. Let $m$ be a multiset of pieces,
$X$ a set of $|m|$ distinct middle blocks of pattern $m$, and $j\in\{1,2\}$.
\begin{enumerate}[label=(\roman*)]
\item For $|m|\in\{1,2\}$ and symmetric parities, the selection of $F_j$ on $X$ has
 the class cost $B-g_F(m)$, and the class of the two aligned matching
 assignments has the class cost $B-g_E(m)$.
\item For $|m|=2$ and skew parities, the aligned fold assignment $(F_j,F_j)$ and
 the set of the two aligned matching assignments satisfy \cref{lem:var-gain}(a)
 with $C=C_m$ and the numbers $b_F(m)$, $b_E(m)$, whose gains are $g_F(m)$ and
 $g^{\rm sel}_E(m)$. \Cref{lem:class-credit} turns them into the credits
 $\kappa_{FF}g_F(m)$ and, on the branch with $\mu$ outside and without background
 or pool terms, $\kappa^{\rm al}_bg^{\rm sel}_E(m)$, against the envelope costs of
 the aligned classes. Crossed matching pairs receive no variance credit.
\item For $|m|\in\{3,4\}$ and an even number of skew parities, not all skew, the
 selection of $F_j$ on $X$ and the class of the two aligned matching assignments
 have the class costs $B-g^{\rm sel}_F(m)$ and $B-g^{\rm sel}_E(m)$. For $|m|=4$
 and four skew parities, the aligned assignments of one kind receive instead,
 inside the envelope class $(4,0)$ or $(0,4)$, the credit
 $\kappa_{4,T}g^{\rm sel}_T(m)$ on their raw-arm prices (\cref{lem:class-credit}).
\item At the per-block upper prices, with $B_*\ge B$ in place of $B$
 (\cref{rem:cp:monotone}) and every repeated-block monomial of a piece multiset
 given the gain of that multiset, these credits lower the complete polynomial of
 \cref{lem:record:record-expansion} by at least $D_{\rm var}$
 \eqref{eq:variance:dvar}, and every coefficient stays nonnegative.
\end{enumerate}
\end{corollary}

\begin{proof}
By the paragraphs above the piece values, the caps and the numbers $\Sigma_i$ are
piece data; the gates of \cref{cert:variance:fresh} give $k_{F,i},k_{E,i}<1$ and
every gain in $[0,1)$; and $R_{\rm err}\le R_*$ by the choice of $R_*$.

(i) \Cref{prop:background-variance}(b) for the fold selection and
\cref{prop:background-variance}(a) for the aligned matching class.

(ii) \Cref{cor:var-classes}(i), (ii) with $C=C_m$ give \cref{lem:var-gain}(a) with
$\bar l_F=b_F(m)$ and $\bar l^{\rm al}_E=b_E(m)$, and $[\bar l_F-R_*]_+=g_F(m)$,
$[\bar l^{\rm al}_E-R_*]_+=g^{\rm sel}_E(m)$. These selections belong to the
envelope classes $(2,0)$ and $(0,2)$ of \cref{cor:envelope-costs}, and the
averaged bound of \cref{lem:var-gain}(a) is the hypothesis of
\cref{lem:class-credit}, whose parts (a) and (b) take the fractions $\kappa_{FF}$
and $\kappa^{\rm al}_b$ of the gains; the background terms of
\cref{prop:background-variance} are not used for this class. The restorations
live on the two aligned corners of the matching bridge
(\cref{lem:record:paired-e-joint}), so no lower bound is available for a crossed
class.

(iii) By \cref{cor:var-classes}(i), (ii) with $C=C_m$ and \cref{lem:var-gain}(c)
the stated classes have the costs $B-g^{\rm sel}_F(m)$, $B-g^{\rm sel}_E(m)$;
they lie in no envelope class, since the classes of \cref{cor:envelope-costs}
consist of skew factors only. For four skew factors the selection lies in the
class $(4,0)$ or $(0,4)$, and \cref{lem:class-credit}(c) credits
$\kappa_{4,T}g^{\rm sel}_T(m)$ on its raw-arm monomials, on its aligned words.

(iv) The per-block upper price of a record of one kind with a parity is half of
the block's share of $tS_{T,i}$ or $tA_{T,i}$, and its raw-arm price half of the
share of $tA^{\rm arm}_{T,i}$. On the monomials of the selections of (i) and (iii)
with a symmetric factor, \cref{lem:record:record-expansion}(iv) allows the
coefficient $B_*-g$, because the class cost is $B-g\le B_*-g$; on the monomials
with a repeated block of the same records, kinds and parities it allows any
coefficient in $[0,B_*c_{\boldsymbol n}]$, and they receive the gain of their
piece multiset, which is below $1<B_*$. On the raw-arm monomials of the classes
of (ii) and of the pure quartics of (iii) the credits are subtracted inside the
replaced classes, as \cref{lem:record:record-expansion}(v) allows, the residual
coefficient being nonnegative by the gates of \cref{lem:class-credit}. It remains
to add up the subtracted amounts.

\emph{Degree one.} The symmetric prices of the two records of one kind of a block
sum, over the blocks of the piece $i$, to $tS_{T,i}$ (the distribution of
\cref{lem:record:record-expansion}(iii) is exact); the gains $g_F(i)$,
$g_E(i)$ give $t\sum_i[g_F(i)S_{F,i}+g_E(i)S_{E,i}]$.

\emph{Degree two.} For a fold record $F_j$, the symmetric monomials of degree two
in its variables form $\frac12(\sum_\Blk x^{\rm s}_{\Blk,F_j})^2$, whose ordered pairs
of blocks carry the price products $\frac{t^2}4S_{F,\Blk}S_{F,\Blk'}$, $S_{F,\Blk}$ being
the share of the block; with the gain of the pair of pieces, summed over $j$ and
over the blocks of each pair of pieces, this is
$\frac{t^2}4\sum_{i,j}g_F(i,j)S_{F,i}S_{F,j}$. The same computation gives
$\frac{t^2}4\sum_{i,j}g_E(i,j)S_{E,i}S_{E,j}$ for the aligned symmetric matching
pairs, the crossed pairs receiving nothing, and, on the raw-arm monomials, the
skew terms $\frac{t^2}4\kappa_{FF}\sum_{i,j}g_F(i,j)A^{\rm arm}_{F,i}A^{\rm arm}_{F,j}$ and
$\frac{t^2}4\kappa^{\rm al}_b\sum_{i,j}g^{\rm sel}_E(i,j)A^{\rm arm}_{E,i}A^{\rm arm}_{E,j}$.

\emph{Degrees three and four.} For a multiset $m$ of size $k$, the monomials of
degree $k$ in the variables of one record of a kind whose blocks lie in the pieces
$m$ are, by the multinomial expansion, those of
$\prod_i(\sum_{\Blk\subseteq[X_i,Y_i)}(x^{\rm s}_\Blk+x^{\rm a}_\Blk))^{m_i}/m_i!$, restricted
to an even number of skew factors; at the per-record prices, and summed over the
two records of one kind, they give $2\cdot2^{-k}t^k\sum_\sigma\prod_lP^{\sigma_l}_{T,i_l}
/\prod_im_i!$ for the words with a symmetric letter, and the same with
$\kappa_{4,T}\prod_lA^{\rm arm}_{T,i_l}$ for the pure quartic. With the gains
$g^{\rm sel}_T(m)$ this is the last line of \eqref{eq:variance:dvar}.

The parts are disjoint: different degrees, kinds, records or parity patterns are
different monomials; the skew credits of degree two and the pure quartic credits
lie inside different envelope classes; and the background and late-pool terms are
inside the gains $g_E(m)$ of the linear and symmetric pair terms only.
\end{proof}

\begin{certificate}[Restorations recomputed at the baseline exponent]\label{cert:variance:fresh}
With the tilted laws on the seven later pieces, the restoration
quantities $M_i$, $\chi_i$ below are formed from the tilted laws'
own marginals (not the uniform ones), all gains lie in $[0,1)$,
$k_{F,i},k_{E,i}<1$, and the assembled total satisfies
$0.00995856031576667<D_{\rm var}<0.00995856031576668$. This is
computed twice, by two independently written assemblies of the
scalar totals in
\path{research/improvements_2026_10_25/global/selected.py}, with
the inputs of \path{research/push_2026_10_25/design.json}. The
adapters in
\path{research/improvements_2026_10_25/global/scalar_adapters.py}
authenticate the exact per-band laws, their stored weights, the
prime thresholds and the single-prime marginal bounds. For each of the seven
changed roots they form, with $n_i=\lfloor570/\log_-X_i\rfloor$,
\[
 M_i=R_{\ulaw}(X_i)^{n_i},\qquad
 \chi_i=(1+d_{\ulaw}(X_i)^2)^{n_i}-1,
\]
rounded outward. The author and independent implementations separately
perform these products. The new marginal is used throughout the
restoration; the fixed uniform half is only the reference measure.
The complete \texttt{new\_root\_restoration} record is in
\path{research/push_2026_10_25/artifacts/author.json}.

The piece product caps are recomputed at anchor $a=\PushA$ from the
exact budget-$570$ products. The selected root changes are already
included in the corresponding native weights. All other piece data,
common-chamber restoration formulas, background credits of degrees
one and two, aligned matching credits, and credits of degrees three
and four are the formulas proved above. Terms pooled over the
pieces from $2^{31}$ are included only when $b_E(m)\ge R_*$.
Both assemblies verify $0\le l\le C_m/2$, every gain in $[0,1)$,
and $k_{F,i},k_{E,i}<1$, then assemble $D_{\rm var}$ using
\eqref{eq:variance:dvar}.

\path{research/push_2026_10_25/verify.py} authenticates the preceding
full native proof and reruns both assemblies with these exact inputs.
Their variance components, product caps, savings at the tail
primes and total closing values are in its two stored outputs;
their
agreement is part of \cref{cert:closing}. In particular the outward
interval for $D_{\rm var}$ lies inside
\[
 0.00995856031576667<D_{\rm var}<0.00995856031576668.
\]
\end{certificate}

\begin{remark}[Sizes]\label{rem:variance:sizes}
At the baseline exponent $a=\PushA$, $D_{\rm var}$ is about $9.96\cdot10^{-3}$. The skew
pair credits dominate: about $6.01\cdot10^{-3}$ for the aligned fold pairs and
$3.88\cdot10^{-3}$ for the aligned matching pairs. The linear fold credit is
about $4.8\cdot10^{-5}$, small because the symmetric parts of the signed fold
records are of degree three; the linear matching credit is about
$1.3\cdot10^{-5}$, the symmetric pairs give less than $6\cdot10^{-10}$, and the
degrees three and four about $5.0\cdot10^{-6}$; the late-pool terms account for
about $1.7\cdot10^{-9}$ of the linear and symmetric-pair credits. Degree six is credited only
through the positive envelopes (\cref{thm:envelopes}).
\end{remark}

\begin{remark}[Robustness]\label{rem:variance:robust}
Every lower bound of this section is an exact relation minus upper bounds of
the kind allowed by (CP8). Step (C) of \cref{lem:record:restoration} subtracts
the energies of the nonroot components from the identity
$Y[\bigotimes_\Blk\mu_\Blk]=Y$; its contraction and density steps, and the
centred interpolations of \cref{lem:var-gain} and \cref{prop:background-variance},
are the exact relations \ref{it:record:P1}--\ref{it:record:P3} and
\ref{it:record:P3b}; the one-sided and bounded-variance steps use the upper bound
$C'$ for normalized energies; the bridges subtract the merge energies; and the
background terms are kept terms of an exact expansion. If every subtracted upper
bound is multiplied by $1+z$, the same proofs give the bounds with $(1+z)C$ in
place of $C$, and the piece caps $(1+z)C_i$; every gain is the larger of $0$ and a
continuous expression, so the strict gates of \cref{cert:variance:fresh} persist
for small $z$ (\cref{lem:record:rec}). The condition $b_E(m)\ge R_*$ of the
late-pool terms holds with room: the strict selected-only margins checked in \cref{cert:variance:fresh}
persist under a sufficiently small perturbation as well.
\end{remark}

\paragraph{Consequences and what to read next.}
Parts II and III: \cref{thm:weil}, \eqref{eq:two-root},
\cref{lem:majorant,def:admissible,lem:sector}, (CP4) and (CP5)
with the charging rule and the monotonicity of the closing
(\cref{def:record:comparison-law,def:cp5,rem:cp:charging,rem:cp:monotone}),
the energy toolkit \ref{it:record:E1}--\ref{it:record:E3},
\ref{it:record:P1}--\ref{it:record:P4} and \ref{it:record:P3b}
(\cref{lem:energy-toolkit}),
\cref{lem:record:variance-interface,lem:record:record-expansion,lem:record:repricing,cor:first-stage-repricing,cor:two-cut-series,lem:record:rec}.
Part IV:
\cref{lem:paley-triple,lem:record:signed-nbhd,sec:frontier:empty-triple-negative,lem:record:projection,lem:matching:kernel,lem:spec:cutkernel,lem:tilt-marginals},
the chamber geometry and polarization
(\cref{lem:record:fold-identification,def:patterns,def:five-type-datum,lem:tr:polarization}).
Appendix A: \cref{thm:ahlgren-ono}.
Part V: the middle blocks, the piece caps, the reference and the bracket
(\cref{def:record:exact-blocks,lem:record:local-domain,lem:record:lower-start-sums,lem:piece-caps,lem:front12:reference,def:record:frame,lem:record:gram}),
the cut L20 of the uniform root (\cref{subsec:middle-l20}, \eqref{eq:l20-cross-matrix}),
the positive roots
(\cref{def:positive-root,lem:positive-root-properties,prop:root-chamber,cert:positive-root-caps}),
the fold records (\cref{lem:old-f-law,lem:record:old-f-marginal,def:record:mixture,lem:signed-components,prop:signed-bounds})
and the matching records (\cref{lem:record:e-law,lem:record:exact-marginal-e}); the
prices (\cref{lem:middle-comparison-prices}); the envelope costs
(\cref{thm:envelopes,lem:class-credit,cor:envelope-costs}); the companion
energies, the sides and the chosen diagonals of the mean
(\cref{lem:record:companion-energy,prop:signed-sides,subsec:matching-envelopes}); the
closing quantities (\cref{def:closing-quantities}); and the data and certificates
\cref{cert:record:weighted-identity,cert:l20-star-mean,cert:variance:fresh,tab:data:piece-values,tab:data:prices,tab:data:positive-roots}.

\paragraph{Hypotheses verified in this section.}
The variance part of (CP5) (\cref{tab:hypotheses}): the classes of
\cref{cor:variance:final}(i),(iii) with their costs, the gains of
\cref{cor:variance:final}(ii) and (iii) that \cref{lem:class-credit} consumes, and
the total $D_{\rm var}$ of \eqref{eq:variance:dvar} that the closing uses; the
diagonal-mean and pool coupons (\cref{lem:record:paired-e-joint}(c),
\cref{lem:pool-coupons}(b)) that \cref{prop:record:mean-credit} uses; and the star
normalizer of the cut L20 (\cref{lem:l20-star-mean}) that
\cref{subsec:middle-l20} uses.
